% =====================================================================
%  Capitolo 12 — Array, timeline e il primo platformer: la seconda
%  lezione di programmazione con i Blueprint di Unreal Engine (Game
%  Design, primo anno)
%  Fonte: trascrizione della lezione GD12 (Game Design, primo anno),
%  seconda lezione di programmazione, diretta prosecuzione del capitolo
%  7. Lezione ibrida: una parte della classe in aula, una parte da
%  remoto, con condivisione dello schermo sul canale vocale "GD1 Voice"
%  di Discord.
%  Contenuti, nell'ordine della lezione: il task assegnato subito (un
%  platformer 3D su tre livelli, da consegnare il 25) e la richiesta di
%  aggiornare il docente come in un team; il PC nuovo che compila gli
%  shader; la domanda sulla compilazione di Unreal dal codice sorgente
%  (modifiche all'engine, server dedicati, console, versioni in
%  anteprima, la 5.9 e la build per browser); Python in Unreal; il
%  peso della build, Unreal Engine 6, i template, Android e Visual
%  Studio; gli array (la classifica, gli indici da zero, Add, Remove,
%  For Each Loop, Get All Actors Of Class, il filtro per distanza);
%  come consegnare (File, Zip Project, le cartelle di un progetto);
%  Maps & Modes; il packaging (Platforms, Shipping, Development,
%  DebugGame); le tre mappe del task; Fab (il launcher, i tipi di
%  contenuto, i filtri, i pacchetti gratuiti a tempo, le versioni
%  dell'engine); la timeline (la piattaforma mobile, l'editor delle
%  curve, il loop, la porta con Play e Reverse, relative e world); il
%  giro tra i banchi (la porta con le box di overlap, il teletrasporto a
%  colpi di Delay, la piattaforma che va avanti e indietro); il
%  movimento sull'Event Tick e i delta seconds; la piattaforma che si
%  distrugge; ricaricare il livello (Get Current Level Name, Open Level
%  by Name); la doppia porta; Engine Scalability; la lava; il timer con
%  i delta seconds; il trampolino (Jump e Launch Character); il mulino
%  e il vento; la scala del personaggio; la chiusura e il riepilogo del
%  task.
%  Integrazioni dal libro "Blueprints Visual Scripting for Unreal
%  Engine 5" (Romero, Sewell), segnalate come approfondimenti.
%
%  Pacchetti richiesti nel preambolo del libro (già presenti nel
%  main.tex): \usepackage{booktabs}, \usepackage{amsmath},
%  \usepackage{graphicx}, \usepackage{tikz} con
%    \usetikzlibrary{shapes.geometric,shapes.symbols,shapes.misc,
%      arrows.meta,positioning,calc,fit}
%  (le librerie vanno caricate nel preambolo, non nel capitolo).
%  Tutte le etichette di questo capitolo hanno il prefisso "gd12:"
%  (gd12:sec:..., gd12:fig:..., gd12:tab:...) e tutti gli stili TikZ
%  "propri" il prefisso "gd12", per non collidere con gli altri
%  capitoli. I grafi Blueprint usano le macro di bp-tikz.tex (la piccola
%  libreria in tasks/gd7-blueprint/), incluse qui sotto così che il
%  capitolo sia autosufficiente; sono protette da un controllo, quindi
%  se un altro capitolo le ha già definite non vengono ridefinite.
%
%  Nota sulla fonte: la trascrizione automatica è buona nelle parti
%  spiegate alla classe (il task, gli array, la consegna, il packaging,
%  Fab, la timeline, il movimento sull'Event Tick, la doppia porta), ma
%  molto rovinata nei lunghi tratti in cui il docente gira tra i banchi
%  o guarda gli schermi condivisi a bassa voce. Quei tratti sono stati
%  trascritti come frasi senza senso o ripetute in loop ("di aziende, di
%  aziende", "mi riguarda, mi riguarda", "piccola piccola piccola",
%  "non ci manchiamo", "lo so, lo so", "non si può fare niente",
%  "è più da vicino", "l'anno solo, ok"), con brani captati da una fonte
%  esterna (frasi su "Signor Presidente, onorevoli colleghi", su un
%  "fascicolo", sul "mercato", su "Sassuolo", sulla "Grecia"): sono
%  artefatti, omessi. Si omettono anche gli scambi privati non legati
%  alla lezione (la nonna che entra nella stanza di uno studente, il
%  bicchiere d'acqua del docente). Diversi nomi sono storpiati ("Arial",
%  "Areal Engine", "il Real", "Arrhenius", "Gigi One Voice", "la Real
%  Tech", "R.I." per "array", "sin" per "scene component"): dove il
%  riferimento è ricostruibile con certezza lo si riporta nella forma
%  corretta; dove non lo è, lo si segnala in nota invece di indovinare.
%  Le figure sono ricostruzioni schematiche dei grafi costruiti a
%  lezione, non screenshot.
% =====================================================================

% --- Macro per i grafi Blueprint (da tasks/gd7-blueprint/bp-tikz.tex) --
% Richiedono TikZ e la libreria calc, già caricati nel preambolo.
\ifcsname bpnode\endcsname\else
\pgfdeclarelayer{bpbg}
\pgfdeclarelayer{bpcomment}
\pgfdeclarelayer{bpwire}
\pgfsetlayers{bpbg,bpcomment,bpwire,main}

% --- colori dell'editor ----------------------------------------------
\definecolor{bpBg}{HTML}{262626}
\definecolor{bpGridFine}{HTML}{2E2E2E}
\definecolor{bpGridBold}{HTML}{1A1A1A}
\definecolor{bpBody}{HTML}{141414}
\definecolor{bpOutline}{HTML}{060606}
\definecolor{bpText}{HTML}{F0F0F0}
\definecolor{bpSub}{HTML}{B4B4B4}
\definecolor{bpValBg}{HTML}{0C0C0C}
\definecolor{bpValBorder}{HTML}{5A5A5A}

% header dei nodi
\definecolor{bpHEvent}{HTML}{8E1B1B}   % eventi (rosso)
\definecolor{bpHFunc}{HTML}{1F5A8C}    % funzioni con esecuzione (blu)
\definecolor{bpHPure}{HTML}{3F7A2E}    % funzioni pure, senza esecuzione (verde)
\definecolor{bpHCast}{HTML}{14706F}    % cast (verde acqua)
\definecolor{bpHFlow}{HTML}{5C5C5C}    % controllo del flusso, macro (grigio)
\definecolor{bpHVar}{HTML}{3A3A3A}     % SET di una variabile

% pin e fili, un colore per tipo di dato (come nell'editor)
\definecolor{bpPinexec}{HTML}{FFFFFF}
\definecolor{bpPinbool}{HTML}{A00000}
\definecolor{bpPinint}{HTML}{1FE3AF}
\definecolor{bpPinfloat}{HTML}{9DF53D}
\definecolor{bpPinstr}{HTML}{F72BD3}
\definecolor{bpPinname}{HTML}{C98AF7}
\definecolor{bpPintext}{HTML}{E2779A}
\definecolor{bpPinvec}{HTML}{F8C624}
\definecolor{bpPinrot}{HTML}{A0B3FC}
\definecolor{bpPintrans}{HTML}{F77A10}
\definecolor{bpPinobj}{HTML}{19A7F2}
\definecolor{bpPinclass}{HTML}{7B2BC4}
\definecolor{bpPinsoft}{HTML}{98F2FF}
\definecolor{bpPinstruct}{HTML}{1458C8}
\definecolor{bpPinwild}{HTML}{8A8A8A}
\definecolor{bpPinenum}{HTML}{0D7A6B}
\definecolor{bpPinskip}{HTML}{141414}

\def\bpRow{0.42}
\def\bpExecStr{exec}
\def\bpSkipStr{skip}
\def\bpCompactStr{compact}

\tikzset{
  bpTitle/.style={text=bpText, font=\sffamily\scriptsize\bfseries,
                  inner sep=0pt, align=left},
  bpTitleC/.style={text=bpSub, font=\sffamily\large\bfseries, inner sep=0pt},
  bpLabel/.style={text=bpText, font=\sffamily\tiny, inner sep=0pt},
  bpNoteSty/.style={text=bpSub, font=\sffamily\tiny\itshape, align=left,
                    inner sep=1pt},
}

\newcommand{\bpval}[1]{{\setlength{\fboxsep}{0.8pt}\setlength{\fboxrule}{0.3pt}%
  \fcolorbox{bpValBorder}{bpValBg}{\textcolor{bpText}{\,#1\,}}}}
\newcommand{\bpsub}[1]{{\color{bpSub}\tiny\itshape\mdseries #1}}

% conta gli elementi non vuoti di una lista "tipo/etichetta, ..."
\newcommand{\bpcountlist}[1]{%
  \gdef\bpres{0}%
  \foreach \bpt/\bpl in {#1}{%
    \ifx\bpt\empty\else
      \pgfmathtruncatemacro\bptmp{\bpres+1}\global\let\bpres\bptmp
    \fi}}

% disegna un pin: tipo, x, y
\newcommand{\bpXpin}[3]{%
  \edef\bpPT{#1}%
  \ifx\bpPT\bpExecStr
    \fill[bpPinexec] (#2-0.085,#3+0.085) -- (#2+0.02,#3+0.085)
       -- (#2+0.105,#3) -- (#2+0.02,#3-0.085) -- (#2-0.085,#3-0.085) -- cycle;
  \else\ifx\bpPT\bpSkipStr
  \else
    \fill[bpPin#1] (#2,#3) circle (0.07);
    \draw[bpPin#1!60!black, line width=0.3pt] (#2,#3) circle (0.07);
  \fi\fi}

% il nodo generico
\newcommand{\bpnode}[8][1]{%
  \bpcountlist{#7}\edef\bpNin{\bpres}%
  \bpcountlist{#8}\edef\bpNout{\bpres}%
  \pgfmathsetmacro\bpRows{max(\bpNin,\bpNout,1)}%
  \edef\bpHC{#4}%
  \ifx\bpHC\bpCompactStr
    \def\bpHead{0.06}%
  \else
    \pgfmathsetmacro\bpHead{0.14+0.34*#1}%
  \fi
  \pgfmathsetmacro\bpH{\bpHead+\bpRows*\bpRow+0.12}%
  \begin{scope}[shift={(#3)}]
    \fill[black, opacity=0.45, rounded corners=4pt]
         (0.07,-0.07) rectangle ++(#6,-\bpH);
    \fill[bpBody, rounded corners=4pt] (0,0) rectangle (#6,-\bpH);
    \ifx\bpHC\bpCompactStr
      \node[bpTitleC] at ({#6-0.52},{-\bpH/2}) {#5};
    \else
      \begin{scope}
        \clip[rounded corners=4pt] (0,0) rectangle (#6,-\bpH);
        \shade[left color=#4, right color=#4!30!bpBody]
              (0,0) rectangle (#6,-\bpHead);
      \end{scope}
      \node[bpTitle, anchor=west] at (0.14,{-\bpHead/2}) {#5};
    \fi
    \draw[bpOutline, line width=0.6pt, rounded corners=4pt]
         (0,0) rectangle (#6,-\bpH);
    \coordinate (#2-c)  at ({#6/2},{-\bpH/2});
    \coordinate (#2-n)  at ({#6/2},0);
    \coordinate (#2-s)  at ({#6/2},-\bpH);
    \coordinate (#2-w)  at (0,{-\bpH/2});
    \coordinate (#2-e)  at (#6,{-\bpH/2});
    \coordinate (#2-sw) at (0,-\bpH);
    \coordinate (#2-ne) at (#6,0);
    \foreach \bpt/\bpl [count=\bpk] in {#7}{%
      \ifx\bpt\empty\else
        \pgfmathsetmacro\bpy{-\bpHead-(\bpk-0.5)*\bpRow-0.06}%
        \bpXpin{\bpt}{0.17}{\bpy}%
        \coordinate (#2-i\bpk) at (0.17,\bpy);
        \node[bpLabel, anchor=west] at (0.33,\bpy) {\bpl};
      \fi}%
    \foreach \bpt/\bpl [count=\bpk] in {#8}{%
      \ifx\bpt\empty\else
        \pgfmathsetmacro\bpy{-\bpHead-(\bpk-0.5)*\bpRow-0.06}%
        \bpXpin{\bpt}{#6-0.17}{\bpy}%
        \coordinate (#2-o\bpk) at (#6-0.17,\bpy);
        \node[bpLabel, anchor=east] at (#6-0.33,\bpy) {\bpl};
      \fi}%
  \end{scope}}

% il GET di una variabile: una "pillola" con un'uscita
\newcommand{\bpget}[5]{%
  \begin{scope}[shift={(#2)}]
    \fill[black, opacity=0.45, rounded corners=6pt]
         (0.06,-0.06) rectangle ++(#5,-0.46);
    \begin{scope}
      \clip[rounded corners=6pt] (0,0) rectangle (#5,-0.46);
      \shade[left color=bpPin#3!40!bpBody, right color=bpBody]
            (0,0) rectangle (#5,-0.46);
    \end{scope}
    \draw[bpOutline, line width=0.6pt, rounded corners=6pt]
         (0,0) rectangle (#5,-0.46);
    \node[bpTitle, anchor=west] at (0.14,-0.23) {#4};
    \bpXpin{#3}{#5-0.17}{-0.23}%
    \coordinate (#1-o1) at (#5-0.17,-0.23);
    \coordinate (#1-n) at ({#5/2},0);
    \coordinate (#1-s) at ({#5/2},-0.46);
  \end{scope}}

% un filo (sotto i nodi)
\newcommand{\bpwire}[3]{%
  \begin{pgfonlayer}{bpwire}
    \edef\bpWT{#1}%
    \ifx\bpWT\bpExecStr\def\bpWW{1.15pt}\else\def\bpWW{0.95pt}\fi
    \draw[bpPin#1, line width=\bpWW]
      let \p1=(#2), \p2=(#3),
          \n1={ifthenelse(\x2<\x1, 45pt, max(12pt, abs(\x2-\x1)*0.45))} in
      (#2) .. controls +(\n1,0) and +(-\n1,0) .. (#3);
  \end{pgfonlayer}}

% riquadro di commento
\newcommand{\bpcomment}[3]{%
  \begin{pgfonlayer}{bpcomment}
    \fill[white, opacity=0.07, rounded corners=3pt] (#1) rectangle (#2);
    \draw[white!45!bpBg, rounded corners=3pt, line width=0.4pt] (#1) rectangle (#2);
    \begin{scope}
      \clip[rounded corners=3pt] (#1) rectangle (#2);
      \path (#1) coordinate (bpcA) (#2) coordinate (bpcB);
      \path ($(bpcA)+(0,-0.42)$) coordinate (bpcH);
      \fill[white, opacity=0.10] (bpcA) rectangle (bpcB |- bpcH);
    \end{scope}
    \node[text=bpText, font=\sffamily\scriptsize\bfseries, anchor=north west,
          inner sep=3pt] at (#1) {#3};
  \end{pgfonlayer}}

\newcommand{\bpnote}[2]{\node[bpNoteSty, anchor=north west] at (#1) {#2};}

% l'ambiente grafo
\newenvironment{bpgraph}[1][]{%
  \begin{tikzpicture}[#1]}{%
  \begin{pgfonlayer}{bpbg}
    \path (current bounding box.south west) ++(-0.35,-0.35) coordinate (bpSW);
    \path (current bounding box.north east) ++(0.35,0.35) coordinate (bpNE);
    \fill[bpBg] (bpSW) rectangle (bpNE);
    \begin{scope}
      \clip (bpSW) rectangle (bpNE);
      \draw[bpGridFine, line width=0.3pt, step=0.25] (bpSW) grid (bpNE);
      \draw[bpGridBold, line width=0.5pt, step=2] (bpSW) grid (bpNE);
    \end{scope}
  \end{pgfonlayer}
  \end{tikzpicture}}
\fi
% ---------------------------------------------------------------------

\chapter[Array, timeline e il primo platformer]{Array, timeline e il primo
platformer: la seconda lezione di programmazione}
\label{gd12:cap:array-timeline}

Il capitolo~\ref{gd7:cap:blueprint} ha aperto la programmazione con i
\textbf{Blueprint}: eventi, nodi e variabili, la trigger box, il \emph{cast},
le monete, i \emph{delta seconds}. Questa è la \textbf{seconda lezione di
programmazione}, e il docente la imposta come la diretta prosecuzione della
prima: ``continuare un attimino sull'impostazione della volta scorsa'', ma su un
\textbf{progetto nuovo}. Come la volta precedente, non sarà una lezione del tipo
``dovete fare questo, questo e questo'': il docente dà \textbf{vari strumenti},
che poi ognuno sfrutta per creare quello che preferisce.

La novità è che, questa volta, il \textbf{task} arriva \emph{all'inizio}: ``vi
assegnerò un task oggi, per comodità ve lo assegno subito, così mentre
sviluppate durante questa lezione potete già cominciare''. Tutto quello che
segue (gli array, la consegna, Fab, la timeline, il movimento, il ricaricamento
del livello) è quindi pensato come \textbf{cassetta degli attrezzi} per quel
task: un piccolo \textbf{platformer 3D su tre livelli}.

La lezione è \textbf{ibrida}: una parte della classe è in aula, una parte
(Mattia, Antonino, e per un tratto Francesco) è collegata da remoto e condivide
lo schermo sul canale vocale di Discord della classe. Anche questa volta il
computer dell'aula dà qualche problema: è stato appena \textbf{cambiato}, e
l'engine, mai aperto prima su quella macchina, deve compilare tutti gli
\emph{shader} prima di partire. La struttura della lezione, e quindi del
capitolo, è questa:
\begin{enumerate}
  \item il \textbf{task}: un platformer 3D su tre livelli, con trappole, aree
        che cambiano il movimento e un livello a tempo, e la richiesta di
        tenere aggiornato il docente ``come in un team''
        (sezione~\ref{gd12:sec:task});
  \item le domande tecniche in attesa che l'engine sia pronto:
        \textbf{compilare Unreal dal codice sorgente}, \textbf{Python},
        Unreal Engine 6, i template, Android e Visual Studio
        (sezioni~\ref{gd12:sec:attesa}--\ref{gd12:sec:altre-domande});
  \item gli \textbf{array}: la classifica, gli indici che partono da zero,
        Add e Remove, il \textbf{For Each Loop}, \textbf{Get All Actors Of
        Class} e il filtro per distanza
        (sezioni~\ref{gd12:sec:array}--\ref{gd12:sec:get-all});
  \item come \textbf{consegnare}: \textbf{Zip Project} e le cartelle di un
        progetto; \textbf{Maps \& Modes}; il \textbf{packaging}
        (sezioni~\ref{gd12:sec:zip}--\ref{gd12:sec:packaging});
  \item le tre mappe del task e \textbf{Fab}
        (sezioni~\ref{gd12:sec:mappe} e~\ref{gd12:sec:fab});
  \item la \textbf{timeline}: la piattaforma mobile, l'editor delle curve, il
        loop, la porta con \emph{Play} e \emph{Reverse}, \emph{relative} e
        \emph{world} (sezioni~\ref{gd12:sec:timeline}--\ref{gd12:sec:relative});
  \item il giro tra i banchi: la porta con le box, il teletrasporto a colpi di
        Delay, la piattaforma che va avanti e indietro
        (sezione~\ref{gd12:sec:giro1});
  \item il \textbf{movimento sull'Event Tick} e i delta seconds, di nuovo
        (sezione~\ref{gd12:sec:tick});
  \item la piattaforma che si distrugge e il \textbf{ricaricamento del
        livello} con \emph{Get Current Level Name} e \emph{Open Level (by
        Name)} (sezioni~\ref{gd12:sec:distrugge} e~\ref{gd12:sec:ricarica});
  \item la \textbf{doppia porta}, costruita dal docente in diretta
        (sezione~\ref{gd12:sec:doppia-porta});
  \item \textbf{Engine Scalability}, la lava, il \textbf{timer}, il
        \textbf{trampolino} (\emph{Jump} e \emph{Launch Character}), il mulino
        e il vento (sezioni~\ref{gd12:sec:scalability}--\ref{gd12:sec:giro2});
  \item la chiusura e il riepilogo del task (sezione~\ref{gd12:sec:chiusura}).
\end{enumerate}

% ---------------------------------------------------------------------
\section{Il task, subito: un platformer su tre livelli}
\label{gd12:sec:task}

\subsection{La consegna}

Il docente manda il testo del task sul canale Discord della classe e condivide
lo schermo per chi è da remoto. ``L'idea sarebbe di creare un
\textbf{platformer 3D} partendo dal \textbf{template di base} di Unreal''; nel
testo ci sono ``tutte quante le informazioni''. In sostanza:

\begin{center}
\fbox{\begin{minipage}{0.92\linewidth}
\vspace{4pt}
\textbf{\large Task per la consegna del 25: un platformer 3D}
\smallskip

\noindent Partendo dal \textbf{template in terza persona} di Unreal Engine,
creare un \textbf{platformer 3D} con questi requisiti minimi:
\begin{enumerate}[leftmargin=1.6em, itemsep=2pt]
  \item \textbf{tre livelli}, cioè tre mappe distinte;
  \item il \textbf{personaggio di base} del template, con telecamera, salto e
        tutto ciò che il template già offre;
  \item degli \textbf{ostacoli particolari}, alcuni dei quali, in certi casi,
        fanno \textbf{ricominciare il livello} (per esempio una trappola che
        uccide il giocatore: il livello riparte);
  \item delle \textbf{aree che modificano il movimento} del giocatore (per
        esempio un'area in cui si va più veloci);
  \item \textbf{un livello con un timer}: se il tempo scade prima che il
        livello sia finito, il livello \textbf{riparte}.
\end{enumerate}
\textbf{Asset:} si possono usare liberamente gli asset di \textbf{Fab}, non
solo quelli del template; \textbf{il codice}, invece, deve essere tutto proprio.
Vanno \textbf{segnalati tutti gli asset presi online}, possibilmente con un
link.\\[2pt]
\textbf{Consegna:} il \textbf{25}, come \textbf{file zip del progetto} (con
\emph{File $\rightarrow$ Zip Project}, sezione~\ref{gd12:sec:zip}) caricato su
\textbf{Google Drive}.\\[2pt]
\textbf{Durante lo sviluppo:} aggiornare il docente sui progressi nella chat
della classe, ``come se fosse una chat di lavoro'' (i messaggi non sono
valutati).
\vspace{4pt}
\end{minipage}}
\end{center}

\paragraph{Le date.} Il progetto si consegna \textbf{il 25}. In realtà la
classe si rivedrà già \textbf{lunedì}, ``però lunedì faremo altre cose''
(la lezione sarà su Git, si veda più avanti). Poi ci si rivede \textbf{il 23}:
il docente lascia tempo fino al 25 ``in modo che, se avete delle domande
all'ultimo, me le fate il giorno stesso''. Nel frattempo, le domande si possono
fare su Discord.

\paragraph{Gli asset da online, e un link.} ``Fondamentalmente dovete
segnalarmi tutti gli asset che utilizzate da online, magari con un link, così
so cosa avete fatto voi e che cosa no.'' Il motivo è ovvio: il task valuta
\textbf{il codice}, e il docente deve poter distinguere il lavoro dello studente
da ciò che è stato scaricato.

\subsection{Il primo e unico task in zip}

Questo sarà ``il \textbf{primo e unico} task che vi darò da consegnare come zip
proprio''. Dai successivi si userà \textbf{GitHub}, ``praticamente un software
per fare in modo che voi possiate caricare man mano che lavorate''; la lezione
di lunedì sarà proprio su quello (``lo vedremo fra quattro, cinque giorni''). Per
questo task, invece: zip, e caricamento su Google Drive.

\subsection{Tenere aggiornato il docente, come in un'azienda}

C'è poi una richiesta nuova, che il docente spiega con cura:
\begin{quote}
Mi piacerebbe un po' come se foste in un'azienda vera, perché in un'azienda vera
è molto difficile che vi diano un task, c'è una data, fatelo. È molto difficile.
\end{quote}
Nella realtà, un task non è ``eccolo, ci vediamo alla scadenza'': chi lo
assegna vuole sapere \textbf{come sta andando}. Per questo il docente chiede di
tenerlo aggiornato mentre si sviluppa, semplicemente scrivendo nella chat della
classe\footnote{La trascrizione dice ``sulla chat di Milano'' (e, all'inizio,
``facciamo Milano''): è con ogni probabilità il canale Discord della classe,
che porta il nome della sede. Non è possibile ricostruire il nome esatto del
canale.} ``come se fosse una chat di lavoro'': ``ho implementato il
\textbf{doppio salto}'', ``ho fatto una \textbf{nuova mappa}'', eccetera, giusto
``per mantenere un progresso''.
\begin{itemize}
  \item Non servono testi formali: ``semplicemente aggiornatelo come se stesse
        lavorando con un team, cosa che poi farete successivamente''.
  \item \textbf{I messaggi non sono valutati}: ``questo è importante''.
\end{itemize}

\paragraph{Approfondimento: lo stand-up, in piccolo.} Quello che il docente
chiede è una versione minima di una pratica comune nei team di sviluppo: il
\emph{daily stand-up} delle metodologie agili, in cui ognuno dice in poche righe
che cosa ha fatto, che cosa farà, che cosa lo blocca (se ne parla nel
seminario su Agile e Scrum, capitolo~\ref{cap:z3-agile-scrum}). Il valore di questi
messaggi non sta nella forma ma nella \textbf{visibilità}: chi coordina vede
presto chi è in difficoltà, e chi lavora si abitua a spezzare il lavoro in
passi piccoli e comunicabili (``il doppio salto'', ``la nuova mappa'').

\subsection{Domande sul task}

Uno studente avvisa che è appena tornato a casa ed è ``al 99\%
dell'installazione di Unreal'', quindi dovrà rifare il progetto da capo:
``tranquillo, perché lo ricominceremo'', anche in classe si parte da un
progetto nuovo. Si può creare subito il progetto per il task, oppure usare
quello della lezione e poi riprenderlo.

\paragraph{Sempre in 3D, sempre in terza persona.} Un'altra domanda: il nuovo
progetto va fatto sempre in 3D? ``Sì, fatelo in 3D, perché tutti quanti i
template di base sono in 3D'': template \textbf{in terza persona}. Più avanti
si vedranno anche il template \textbf{in prima persona} e quello del
\textbf{veicolo}, ``li vedremo un pochettino tutti tranne quello VR, purtroppo'',
e quello AR, ``che tanto non serve particolarmente, non si usa più tanto''.

% ---------------------------------------------------------------------
\section{In attesa degli shader}
\label{gd12:sec:attesa}

\subsection{Il PC nuovo, e GD1 Voice}

Il computer dell'aula è appena stato cambiato, e l'engine ``non è mai stato
caricato su questa versione'': deve compilare gli \emph{shader}, e il launcher
(``si è un po' buggato'') non vuole scaricare i pacchetti. ``Speriamo di no,
però\dots{} non capisco perché non mi ha messo template e funzionalità.''

Nell'attesa, il docente chiede a chi è da remoto di entrare nel canale vocale
\textbf{GD1 Voice}\footnote{Trascritto anche come ``Gigi One Voice''.} di
Discord e di \textbf{condividere lo schermo}; chi è in aula può anche non farlo,
``tanto lo vedo''. E chiede a tutti di \textbf{creare un nuovo progetto in
terza persona di base}, così da poter dare una mano mentre si lavora. Uno
studente avvisa che Unreal gli va già a scatti, e con la condivisione su
Discord ``esplode il computer'': ``va bene, condividi quando ti serve''.

\subsection{Che cosa avete fatto dalla volta scorsa?}

Il docente chiede se ci sia qualcosa della lezione precedente che non è
riuscito, o se qualcuno abbia scritto codice per altri progettini personali. Uno
studente ha fatto qualcosa con una \textbf{timeline} (``timeline che non vi ho
spiegato la volta scorsa, quindi oggi ve la spiego sicuramente'', e infatti è
il cuore della lezione, sezione~\ref{gd12:sec:timeline}).

\paragraph{L'editor in inglese, di nuovo.} Quando finalmente l'engine si apre,
è in italiano: ``non posso sopportarlo, devo assolutamente toglierlo''. Il
docente lo rimette in inglese dalle preferenze dell'editor (\emph{Region \&
Language}), come già spiegato nel capitolo~\ref{gd7:cap:blueprint}
(sezione~\ref{gd7:sec:inglese}), e toglie anche l'overlay di Discord, ``veramente
fastidioso''.

% ---------------------------------------------------------------------
\section{Compilare Unreal dal codice sorgente}
\label{gd12:sec:sorgente}

\subsection{La domanda}

Uno studente remoto ha l'engine che si sta ancora caricando: sta
\textbf{compilando Unreal da zero}. Un compagno chiede: che cosa cambia tra
compilare il codice sorgente e usare la versione di Unreal ``già fatta''?
``Domanda molto interessante.''

\subsection{Un engine con il codice aperto}

Premessa: Unreal Engine si può semplicemente scaricare dall'\textbf{Epic Games
Launcher}, si preme \emph{Play} ``e siamo tutti quanti a posto e contenti''. Ma
Unreal è uno dei pochi engine (a differenza, per esempio, di Unity) che
\textbf{mette a disposizione il codice sorgente}: passando da
GitHub\footnote{La trascrizione dice ``andando su di tab'': il codice sorgente
di Unreal è distribuito su GitHub, dopo aver collegato il proprio account Epic
a quello GitHub.} si può accedere al codice di tutte le versioni,
scaricarlo e \textbf{compilarlo da sé}. Ci sono vantaggi e svantaggi.

\paragraph{Lo svantaggio.} ``Pesa un sacco, e ci vuole un giorno per fare la
build, fondamentalmente.'' Uno studente ipotizza ``tipo 800 giga''; il docente
non conferma la cifra, ma ricorda che una volta che l'aveva fatta completa
``mi pesava un sacco'', e sospetta che allo studente manchi qualcosa.

\paragraph{I vantaggi.} I casi in cui conviene sono, nella sostanza, quattro.
\begin{enumerate}
  \item \textbf{Modificare l'engine.} Se si sa programmare in C++, si può
        cambiare l'engine come si vuole. Serve quando manca una funzionalità
        che \emph{non si può mettere neanche come plugin}: ``volete che
        tutte le schermate dell'engine, quando le muovete, siano
        gommose?''. Non si può fare con un plugin, si può fare così. ``La
        maggior parte delle funzionalità si possono mettere con plugin, quindi
        non serve'', ma in quei casi sì.
  \item \textbf{I server dei giochi multiplayer.} Per un gioco multiplayer
        serve la build \emph{non per il computer, ma per il server} (il
        \emph{dedicated server}), che con l'Unreal ``di base'' non si può
        fare.
  \item \textbf{Le console.} Per buildare per PlayStation, Nintendo e simili
        servono dei file che forniscono Sony, Nintendo e gli altri: bisogna
        compilare l'engine \emph{con quei file dentro}, e solo allora si può
        fare la build per quelle piattaforme.
  \item \textbf{Le versioni in anteprima.} Si possono scaricare le versioni
        ``super preview'', ancora prima che vengano rilasciate. Su GitHub c'è
        già anche \textbf{Unreal Engine 6}: il docente la sta compilando
        proprio per testarne le funzioni, e anticipa che, quando uscirà, ci
        sarà ``sicuramente'' un corso di aggiornamento. ``Comunque i concetti
        di base, float eccetera, sono tutti quanti in tutti i linguaggi di
        programmazione. Quindi non vi preoccupate.''
\end{enumerate}
``Non credo che vi servirà a breve'', conclude, ma può essere utile saperlo.

\subsection{La 5.9 e la build per browser}

Nel launcher, nell'elenco delle versioni, in fondo c'è già la \textbf{5.9}
(``l'hanno messa''). È ``uguale alle altre'', se non per una feature ``molto
carina'': si può fare la \textbf{build per browser}. Se si vuole pubblicare il
proprio gioco su piattaforme come \textbf{itch.io}, giocabile nel browser
``senza che uno debba scaricarlo'', ``questo è ottimo''. Era una cosa che
c'era già in Unreal Engine~4 (il docente cita la 4.22), poi l'hanno tolta, e
ora l'hanno rimessa.\footnote{Il supporto HTML5 di Unreal Engine~4 è stato
effettivamente rimosso dalla distribuzione ufficiale intorno alla 4.24 e reso
un progetto della comunità. Della build per browser della 5.9 si riporta solo
quanto detto dal docente, che non l'ha ancora provata.} Avvertenza: è una
versione ``super preview'', ``non l'ho ancora provata, probabilmente è
terribile, però si fa quel che si può''.

% ---------------------------------------------------------------------
\section{Altre domande su Unreal}
\label{gd12:sec:altre-domande}

\subsection{Python in Unreal}

Un'altra domanda riguarda il motivo per cui lo studente si sta compilando
l'engine: ci sono altri linguaggi oltre ai Blueprint e al C++? Sì: in Unreal si
possono creare \textbf{tool in Python}; nella sezione dei \textbf{plugin}
c'è proprio il plugin \emph{Python}. Il docente precisa che cosa ci si fa:
\begin{itemize}
  \item \textbf{non il codice di gioco}, ``perché sarebbe anche un pochettino
        pesante, essendo che Python non è proprio real time'';
  \item ma \textbf{script per l'editor}, per esempio per agire su tutti gli
        asset contemporaneamente.
\end{itemize}
Queste operazioni sull'editor si possono fare con i Blueprint, con il C++ e con
Python: ``è tutto tra le possibilità, quindi fai come vuoi''.

\paragraph{L'esempio del docente.} In un gioco a cui ha lavorato, aveva creato
un tool che, \textbf{premendo un tasto}, comprimeva tutte le texture in modo
diverso \textbf{in base alla cartella} in cui si trovavano: quelle dei
personaggi con una certa qualità, quelle delle mappe con un'altra, e così via,
``in base a vari parametri''. ``Può essere utile, non è indispensabile.'' (Lo
studente, comunque, non si sta compilando l'engine per usare Python: non c'è
bisogno di farlo per quello.)

\subsection{Android e Visual Studio}

A uno studente ``un pochettino smanettone'' il docente fa notare che la build
custom potrebbe tornare comoda anche per \textbf{Android}, ``ma in realtà si può
fare anche senza'': bisogna però installare ``tipo un milione di software di
Android'', Android Studio eccetera.

Più avanti nel percorso si vedrà anche come installare \textbf{Visual Studio}
con tutto il necessario. Finora si sono fatti progetti \textbf{full Blueprint},
``quindi senza codice vero''. Ma se si entra in un team con un progetto
\textbf{metà Blueprint e metà C++}, Visual Studio serve, ``se no non potete né
aprirlo né fare niente''; lo stesso vale se il progetto ha plugin o ``robe
particolari'' non già compilate. In quel momento, anticipa il docente, si
scaricherà direttamente la versione successiva, perché ci sarà Unreal Engine~6.

% ---------------------------------------------------------------------
\section{Gli array}
\label{gd12:sec:array}

\subsection{Una variabile, tante informazioni}

Con l'engine finalmente pronto (``vi annuncio subito di che cosa cominceremo a
parlare''), il primo argomento sono gli \textbf{array}: ``delle variabili, però
un po' particolari''. Le variabili del capitolo~\ref{gd7:cap:blueprint}
(sezione~\ref{gd7:sec:variabili}), Boolean, Float, Integer e così via,
mantengono \textbf{un'informazione} ciascuna. Una funzionalità ``molto
importante'' è poter contenere, \textbf{in una variabile singola, più
informazioni dello stesso tipo}: questo è un array.

\paragraph{L'esempio della classifica.} Si deve fare una \textbf{leaderboard},
con diversi punteggi. Senza array, si potrebbe creare una variabile
\emph{Primo} (un intero con il punteggio di chi arriva primo), poi una
\emph{Secondo}, poi una \emph{Terzo}, e così via. ``Chiaramente diventa molto
difficile poi da gestire il codice'': con cinquanta posizioni ``diventerebbe un
casino''. Con un array, invece, si ha una sola variabile che contiene tutti i
punteggi, in ordine.

\subsection{Creare un array}

Si crea una variabile normale e, accanto al suo tipo, si clicca sul
\textbf{pallino} (l'icona del contenitore) e si sceglie \textbf{Array}. In
questo modo l'informazione singola diventa un array, che ``permette di
incasellare infinite, tra virgolette, informazioni dello stesso tipo''.
\begin{enumerate}
  \item Si \textbf{compila} il Blueprint (come per tutte le variabili, è il
        momento in cui compare il pannello del valore di default).
  \item Nel valore di default si preme \textbf{+} per aggiungere elementi:
        il primo della classifica ha \textbf{950} punti, il secondo
        \textbf{800}, il terzo \textbf{700}, ``e così via. E ne posso mettere
        quanti ne voglio.''
\end{enumerate}
Nel grafo, un array si riconosce dall'icona: al posto del pallino del tipo
singolo c'è una \textbf{griglia di quadratini}\footnote{A voce il docente dice
``quando ci sono questi 8 pallini, questi 8 quadratini''; l'icona dell'array
nell'editor è una piccola griglia $3 \times 3$.}. ``Vuol dire che è un array,
perché altrimenti si vedrebbe così'', con il pallino normale.

\subsection{Gli indici partono da zero}

``Informazione importante sono gli \textbf{indici}.'' I valori di un array
sono \textbf{numerati}, come ci si aspetta, \emph{ma la numerazione parte da
0}. Il primo valore di un array ha indice \textbf{0}, il secondo indice 1, poi
2, 3, 4 e così via.

\begin{table}[htbp]
\centering
\begin{tabular}{lccccc}
\toprule
\textbf{Posizione in classifica} & 1\textsuperscript{a} & 2\textsuperscript{a}
  & 3\textsuperscript{a} & 4\textsuperscript{a} & \dots \\
\midrule
\textbf{Indice nell'array} & 0 & 1 & 2 & 3 & \dots \\
\textbf{Valore} (Leaderboard) & 950 & 800 & 700 & 8 & \dots \\
\bottomrule
\end{tabular}
\caption{L'array \emph{Leaderboard} della lezione, dopo un \emph{Add} di 8: il
primo elemento ha indice 0, l'ultimo ha indice pari alla lunghezza meno uno.}
\label{gd12:tab:indici}
\end{table}

\paragraph{Non solo numeri.} Un array è utile ``non solo con i numeri'', anche
con gli \textbf{oggetti}: per esempio un array di \textbf{tutte le monete nella
scena}, così che con un tasto si prendano tutte quelle intorno. In generale,
serve ovunque ci sia bisogno di più di un'informazione dello stesso tipo:
punteggi, oggetti, posizioni.

\paragraph{Una cattiva abitudine da evitare.} ``Cercate di evitare di fare, in
qualsiasi caso, variabili diverse se sono la stessa cosa soltanto in posizioni
diverse. Vedo molto spesso fare tanti interi singoli al posto di fare un array,
semplicemente.''

\subsection{Add e Remove}

Come si usano? Il docente aggancia il codice a un tasto qualsiasi, la
\textbf{E} (``giusto per vedere come scatenare gli eventi''), e trascina nel
grafo il \textbf{GET} dell'array. Dal suo pin si possono tirare fuori i nodi
propri degli array.
\begin{itemize}
  \item \textbf{Add} aggiunge un valore nuovo \textbf{in fondo} alla lista. Con
        i primi tre valori (950, 800, 700), un Add di 8 dà ``questa cosa qua
        più 8, qui'': 950, 800, 700, 8.
  \item \textbf{Remove} (nell'editor: \emph{Remove Index}) toglie un elemento:
        qui bisogna dirgli \textbf{quale posizione} rimuovere. L'indice 0
        ``vi toglierà il primo dalla classifica, ve lo toglie semplicemente''.
\end{itemize}

\paragraph{Approfondimento: gli altri nodi degli array.} Il libro di Romero e
Sewell elenca i nodi più usati, che conviene conoscere da subito:
\textbf{Get (a copy)} legge l'elemento a un certo indice; \textbf{Set Array
Elem} lo sovrascrive; \textbf{Insert} aggiunge un elemento \emph{a un indice
dato}, spostando avanti tutti i successivi (a differenza di Add, che aggiunge in
fondo); \textbf{Length} restituisce il numero di elementi e \textbf{Last Index}
l'indice dell'ultimo, che è sempre \emph{Length}~$-$~1; \textbf{Contains} dice
se un valore è presente; \textbf{Random Array Item} pesca un elemento a caso;
\textbf{Is Empty} e \textbf{Is Not Empty} controllano se l'array ha elementi. Il
libro avverte anche di \textbf{non leggere mai un indice oltre l'ultimo}: è un
errore che può dare risultati imprevisti e crash difficili da rintracciare.

\subsection{Il For Each Loop: vederli tutti}

``La cosa più utile che si fa con gli array non è aggiungere e togliere, ma è
\textbf{lopparli}, ovvero vederli tutti quanti.'' Per stampare la lista delle
posizioni, per esempio, si fa un \textbf{loop}: trascinando dal pin dell'array
si sceglie \textbf{For Each Loop}, ``un nodo un pochettino complesso rispetto
agli altri, ma neanche troppo''. Per non confondere gli interi del punteggio con
quelli dell'indice, il docente cambia il tipo dell'array in \textbf{Float} e lo
rinomina \textbf{Leaderboard}.

Il For Each Loop ha quattro uscite.
\begin{itemize}
  \item \textbf{Loop Body}: un'uscita di esecuzione che parte \textbf{una
        volta per ogni elemento} (``qua mi esce otto volte l'output, rispetto a
        quanti elementi ci sono qua dentro'').
  \item \textbf{Array Element}: il valore dell'elemento corrente: ``mi esce il
        primo valore, poi il secondo, poi il terzo, finché non finisce''.
  \item \textbf{Array Index}: l'indice dell'elemento corrente: 0, 1, 2, 3, 4,
        ``finché non finisce''.
  \item \textbf{Completed}: parte una volta sola, quando il ciclo è finito.
\end{itemize}

\begin{figure}[htbp]
\centering
\begin{bpgraph}
  \bpcomment{-0.3,0.75}{12.6,-3.0}{Stampare la classifica con il tasto E}
  \bpnode{e}{0,0}{bpHEvent}{E}{2.2}{}{exec/Pressed, exec/Released, struct/Key}
  \bpnode{fe}{3.3,0}{bpHFlow}{For Each Loop}{3.4}{exec/Exec, float/Array}%
    {exec/Loop Body, float/Array Element, int/Array Index, exec/Completed}
  \bpget{lb}{0.4,-1.9}{float}{Leaderboard}{2.2}
  \bpnode{cv}{7.5,-1.45}{compact}{}{0.8}{float/}{str/}
  \bpnode{ps}{9.2,0}{bpHFunc}{Print String}{3.0}{exec/, str/In String}{exec/}
  \bpwire{exec}{e-o1}{fe-i1}
  \bpwire{float}{lb-o1}{fe-i2}
  \bpwire{exec}{fe-o1}{ps-i1}
  \bpwire{float}{fe-o2}{cv-i1}
  \bpwire{str}{cv-o1}{ps-i2}
  \bpnote{3.0,-2.35}{\emph{Leaderboard} è un array di float (950, 800, 700, \dots)}
\end{bpgraph}
\caption{Il primo uso di un array: alla pressione di E, il For Each Loop esegue
il Print String una volta per ogni elemento della classifica. Sullo schermo i
valori compaiono dal basso verso l'alto, perché l'ultimo messaggio va in cima.}
\label{gd12:fig:foreach}
\end{figure}

Collegato un \textbf{Print String} al Loop Body, con l'Array Element come
testo, si preme Play e poi E: ``vedete che mi escono in ordine''. Anzi, sullo
schermo sembrano al contrario, ``semplicemente perché i messaggi dopo vanno in
cima'': ogni nuovo messaggio spinge giù i precedenti. ``Tuttavia è corretto.''

% ---------------------------------------------------------------------
\section{Get All Actors Of Class e il filtro per distanza}
\label{gd12:sec:get-all}

\subsection{Prendere tutti gli oggetti di un tipo}

Il For Each Loop diventa davvero potente con un nodo che \emph{restituisce} un
array: \textbf{Get All Actors Of Class}, ``prendi tutti quanti gli oggetti di un
determinato tipo''. Il docente crea al volo un Blueprint Class
\textbf{Coin}, ci mette dentro ``una bella sfera'', ne sparge un po' nel
livello, e torna sul codice del personaggio:
\begin{enumerate}
  \item alla pressione di E, \textbf{Get All Actors Of Class} con classe
        \emph{Coin}: ``prendimi tutte quante le coin'';
  \item dal suo array di uscita, un \textbf{For Each Loop};
  \item nel Loop Body, un \textbf{Destroy Actor} con target l'Array Element.
\end{enumerate}
Play, E: le monete \textbf{scompaiono tutte}. ``Concettualmente vi torna? Tutto
quanto chiaro? Chiaramente dovrete prenderci un attimino la mano.''

\subsection{Controllarli uno per uno: la distanza}

``La cosa comoda di avere questo loop è controllarli \textbf{singolarmente}.''
Per esempio: prendere solo le monete \textbf{entro un certo raggio} dal
personaggio.
\begin{itemize}
  \item Dall'Array Element si prende la sua posizione, \textbf{Get Actor
        Location} (``come abbiamo già fatto la volta scorsa''): è la posizione
        della moneta.
  \item Un altro \textbf{Get Actor Location}, senza target, prende
        \emph{self}, cioè il personaggio.
  \item Il nodo \textbf{Distance} (\emph{Vector Distance}) restituisce la
        distanza tra i due punti.
  \item Se la distanza è \textbf{minore di 200}, un \textbf{Branch} manda il
        flusso al Destroy Actor; altrimenti no.
\end{itemize}

\begin{figure}[htbp]
\centering
\begin{bpgraph}
  \bpnode{e}{0,0}{bpHEvent}{E}{2.0}{}{exec/Pressed, exec/Released, struct/Key}
  \bpnode[2]{ga}{2.8,0}{bpHFunc}{Get All Actors Of Class\\\bpsub{Actor Class: Coin}}{3.6}%
    {exec/, class/Actor Class \bpval{Coin}}{exec/, obj/Out Actors}
  \bpnode{fe}{7.3,0}{bpHFlow}{For Each Loop}{3.2}{exec/Exec, obj/Array}%
    {exec/Loop Body, obj/Array Element, int/Array Index, exec/Completed}
  \bpnode{br}{11.4,0}{bpHFlow}{Branch}{2.4}{exec/, bool/Condition}{exec/True, exec/False}
  \bpnode[2]{da}{12.0,-2.4}{bpHFunc}{Destroy Actor\\\bpsub{Target is Actor}}{2.6}%
    {exec/, obj/Target}{exec/}
  \bpnode{loc1}{3.0,-3.2}{bpHPure}{Get Actor Location}{3.3}{obj/Target}{vec/Return Value}
  \bpnode{loc2}{3.0,-4.6}{bpHPure}{Get Actor Location}{3.3}{obj/Target \bpval{self}}{vec/Return Value}
  \bpnode{dist}{7.0,-3.6}{bpHPure}{Distance}{2.2}{vec/V1, vec/V2}{float/}
  \bpnode{lt}{9.6,-3.6}{compact}{$<$}{1.7}{float/, float/\bpval{500}}{bool/}
  \bpwire{exec}{e-o1}{ga-i1}
  \bpwire{exec}{ga-o1}{fe-i1}
  \bpwire{obj}{ga-o2}{fe-i2}
  \bpwire{exec}{fe-o1}{br-i1}
  \bpwire{exec}{br-o1}{da-i1}
  \bpwire{obj}{fe-o2}{loc1-i1}
  \bpwire{obj}{fe-o2}{da-i2}
  \bpwire{vec}{loc1-o1}{dist-i1}
  \bpwire{vec}{loc2-o1}{dist-i2}
  \bpwire{float}{dist-o1}{lt-i1}
  \bpwire{bool}{lt-o1}{br-i2}
\end{bpgraph}
\caption{Distruggere solo le monete vicine: per ogni moneta della scena si
calcola la distanza dal personaggio, e un Branch lascia passare solo quelle
entro 500 unità (5 metri). Lo schema ``prendi tutti, controlla uno per uno,
agisci'' è alla base di moltissimo codice.}
\label{gd12:fig:distanza}
\end{figure}

Il valore va tarato a occhio, provando: 200 ``è un pochettino poco'', 2000
prende ``tutte quante quelle di quest'area'' ed è troppo; con \textbf{500}
``tac, tutte quante quelle nell'area''. Il principio è importante:
\begin{quote}
Con questo concetto potete veramente sbizzarrirvi. Diciamo che buona parte del
codice viene fatto con questo principio qua. Ed è super importante.
\end{quote}
Domande sugli array? Nessuna, per ora: ``chiaramente vi verranno dei dubbi
dopo, non vi preoccupate''.

\paragraph{Approfondimento: un nodo costoso.} Get All Actors Of Class scorre
\emph{tutti} gli attori del livello per trovare quelli della classe richiesta:
va benissimo alla pressione di un tasto, come qui, o una volta al Begin Play,
ma è una cattiva idea chiamarlo \textbf{a ogni frame} sull'Event Tick in un
livello pieno di oggetti. Se serve spesso, conviene chiamarlo una volta e
salvare il risultato in una variabile array.

% ---------------------------------------------------------------------
\section{Come consegnare: Zip Project}
\label{gd12:sec:zip}

``Allora, questo è importante per la consegna.'' Anche quando si userà Git, il
docente vorrà sempre \textbf{poter vedere il progetto}, non solo
l'eseguibile.\footnote{La trascrizione dice ``non solamente la VIP'':
plausibilmente la \emph{build}, cioè il gioco esportato.} Come si condivide un
progetto ``senza utilizzare metodi strani, GitHub eccetera''?
\begin{itemize}
  \item \textbf{Sì:} \emph{File $\rightarrow$ \textbf{Zip Project}}, si
        sceglie il nome, si salva, e dopo qualche secondo lo zip è pronto.
  \item \textbf{No:} selezionare la cartella del progetto in Esplora risorse e
        comprimerla a mano, ``perché sennò vi metto un sacco di cartelle che
        non servono e pesano un sacco''.
\end{itemize}

\subsection{Che cosa c'è in una cartella di progetto}

Per capire il perché, il docente apre la cartella di un progetto Unreal.

\begin{table}[htbp]
\centering
\begin{tabular}{p{3.2cm}p{9.8cm}}
\toprule
\textbf{Cartella} & \textbf{Che cosa contiene, e se serve} \\
\midrule
\texttt{Config} & Le impostazioni del progetto: ``molto importante''. \\
\texttt{Content} & Tutti gli asset: mappe, Blueprint, mesh, materiali,
  texture. ``La cosa più importante in assoluto.'' \\
\texttt{Binaries}, \texttt{DerivedDataCache}, \texttt{Intermediate} & File
  generati dall'engine e cache: ``non servono a niente'' per la consegna,
  vengono \textbf{rigenerati} ogni volta che si apre il progetto. \\
\texttt{Saved} & Log, autosalvataggi e \textbf{salvataggi di gioco}: cancellarla
  azzera anche i salvataggi. \\
file \texttt{.uproject} & Il file del progetto, quello che si apre con
  il doppio clic. \\
\bottomrule
\end{tabular}
\caption{Le cartelle tipiche di un progetto Unreal. Zip Project tiene ciò che
serve ed esclude le cartelle rigenerabili (e quelle in più che compaiono nei
progetti C++).}
\label{gd12:tab:cartelle}
\end{table}

Zip Project esclude da solo le cartelle che si possono rigenerare, e le altre
che compaiono nei progetti in C++, ``quindi ci pensa lui senza che dobbiate farlo
voi''.\footnote{La trascrizione elenca a voce ``queste tre qua non servono a
niente\dots{} questa è cache, questa anche. Qua dentro mi tiene i log e i
salvataggi di gioco'' senza nominarle; la tabella riporta la struttura standard
di un progetto Unreal.}

% ---------------------------------------------------------------------
\section{Maps \& Modes: quale mappa si apre}
\label{gd12:sec:maps-modes}

Il task chiede più mappe, e quindi bisogna sapere con chiarezza \textbf{quale
mappa si apre}: quella che si apre allo sviluppatore quando apre l'engine
(``è quello meno importante''), ma soprattutto quella che si apre \textbf{al
giocatore} quando avvia il gioco. Si imposta da \emph{Edit $\rightarrow$
Project Settings $\rightarrow$ \textbf{Maps \& Modes}}, dove ci sono due campi:
\begin{itemize}
  \item \textbf{Editor Startup Map}: la mappa che si apre allo sviluppatore
        quando apre il progetto nell'editor;
  \item \textbf{Game Default Map}: la mappa che si apre al giocatore quando
        avvia il gioco.
\end{itemize}
Il gioco si aprirà \textbf{sempre} su questa mappa, che di solito è il
\textbf{main menu}, oppure una \textbf{schermata di avvio} con il logo dello
studio, che a sua volta apre il main menu: ``tendenzialmente funziona così''.

% ---------------------------------------------------------------------
\section{Il packaging: dare il gioco a chi non ha l'engine}
\label{gd12:sec:packaging}

Fare la \textbf{build} non è richiesto per il task, ``più che altro perché ogni
tanto su alcuni PC può dare problemi, ci può volere un po', eccetera''. Ma il
docente mostra come si fa a far giocare chiunque, anche chi non ha l'engine.

\paragraph{Platforms.} Nella barra in alto c'è il menu \textbf{Platforms}, con
la lista di tutti i dispositivi per cui si può buildare. Sul PC dell'aula sono
tutti con un \textbf{punto esclamativo}: mancano i componenti per fare davvero
la build per quelle piattaforme (Android, Linux eccetera, ``bisogna attivare un
po' di cose''). Chi compila Unreal da sé con i file delle console
(sezione~\ref{gd12:sec:sorgente}) se le ritrova qui, tra le altre.

\paragraph{Le configurazioni.} Prima di fare \emph{Package Project} si sceglie
la configurazione della build, ``abbastanza intuitive, non particolarmente
complesse''.

\begin{table}[htbp]
\centering
\begin{tabular}{p{2.6cm}p{10.4cm}}
\toprule
\textbf{Configurazione} & \textbf{A che cosa serve, secondo il docente} \\
\midrule
\textbf{Shipping} & La versione finale, ``quella proprio che potete dare
  pubblica a tutti quanti''. Toglie tutti gli strumenti di debug, perché si
  presuppone che la si pubblichi (su Steam e simili). \\
\textbf{Development} & ``Abbastanza inutile'', nell'uso del docente. \\
\textbf{DebugGame} & La più utile mentre si sviluppa: permette di usare la
  \textbf{console dei comandi}, vedere i \textbf{log} del gioco, fare
  \textbf{test} (per esempio sul lag) più facilmente. \\
\bottomrule
\end{tabular}
\caption{Le configurazioni di packaging. In termini tecnici anche Development
conserva console e log; DebugGame in più compila il codice del gioco senza
ottimizzazioni, per poterlo seguire con il debugger.}
\label{gd12:tab:packaging}
\end{table}

\paragraph{La procedura.} Si sceglie \textbf{Shipping}, poi \emph{Package
Project}. L'editor avvisa che manca un SDK: si fa \emph{Continua}. Si sceglie
una cartella di destinazione: il docente di solito crea una cartella
\texttt{Package} dentro il progetto, ``però poi a vostra discrezione''. ``Ci
vorrà una vita, una vita'', soprattutto la prima volta; le volte successive
meno, perché l'engine ``si ricorda'' parecchio di ciò che ha già elaborato. (In
classe il docente interrompe il packaging appena partito.)

% ---------------------------------------------------------------------
\section{Le tre mappe del task}
\label{gd12:sec:mappe}

``Allora io direi di cominciare proprio con il task.'' Il primo passo è creare
\textbf{le tre mappe}. Il docente ricorda che \textbf{il task sono i requisiti
minimi}: ``tutto quanto ciò che ci sta sopra lo potete creare tranquillamente''.
La \textbf{gerarchia delle cartelle} non è importante, ``strutturatele un po'
come preferite''; il docente crea per esempio una cartella \texttt{Maps}.
\begin{enumerate}
  \item Non creare il livello con clic destro $\rightarrow$ \emph{Level} nel
        Content Browser, ``perché sennò ve lo crea veramente vuoto'' (si può
        anche fare, ma poi manca tutto).
  \item Fare invece \emph{File $\rightarrow$ \textbf{New Level}} e partire
        almeno dal \textbf{Basic}, ``così avete almeno il cielo e le cose di
        base''. (Su come mettere le luci in scena, il docente rimanda al corso
        di environment; si veda anche il capitolo~\ref{gd11:cap:blockout-luci}.)
  \item Mettere in scena un oggetto e \textbf{salvare}: \emph{Level1}.
  \item Per le altre due, \textbf{Ctrl+C, Ctrl+V} sul livello nel Content
        Browser: \emph{Level2} e \emph{Level3}. ``Perfetto, fantastico.''
\end{enumerate}
Sono i requisiti minimi: ``se ne volete fare di più, fatene tranquillamente di
più''.

% ---------------------------------------------------------------------
\section{Fab: gli asset da online}
\label{gd12:sec:fab}

\subsection{Dove si trova}

In altri task gli studenti avevano restrizioni che non permettevano di usare
\textbf{Fab}, il negozio di asset di Epic. Qui invece sì, e il docente spiega
come si usa. Dipende un po' dalla versione dell'Epic Games Launcher, ``vecchio o
nuovo, è abbastanza indifferente, però alcune cose potrebbero cambiare'':
\begin{itemize}
  \item nel launcher nuovo c'è una categoria \textbf{Fab};
  \item nel vecchio, sotto la lista degli engine e dei progetti, c'è il
        \textbf{Vault} (la libreria), e da lì si apre Fab: ``un pochettino più
        macchinoso''.
\end{itemize}
Fab è disponibile anche \textbf{dentro l'editor}, dal menu \emph{Window}
$\rightarrow$ \emph{Fab}, e da lì si possono aggiungere i contenuti direttamente
al progetto.

\subsection{Asset sì, codice no}

``Non vi do restrizioni sugli asset, chiaramente. Non codice, ok?'' Il codice
deve essere \textbf{tutto proprio}, ``perché sennò è un po' cheattare''.
Per il resto, ci sono ``un sacco di pacchetti gratuiti''. E i pacchetti
piratati? ``No, no, no, e se l'hanno fatto è illegale.''

\subsection{Tre tipi di contenuto}

Su Fab i contenuti sono di tre tipi.
\begin{enumerate}
  \item \textbf{Asset}: si fa \emph{Add} e vengono aggiunti a un progetto.
  \item \textbf{Plugin}: si installano \emph{sulla versione dell'engine}, non
        in un progetto specifico (``anche se adesso questa opzione l'hanno un
        pochettino cambiata'').
  \item \textbf{Progetti completi}: si fa \emph{Crea progetto} e viene creato
        un progetto da zero con quel contenuto.
\end{enumerate}

\subsection{Filtrare bene}

Fab non contiene solo roba per Unreal: ``prima c'era il Marketplace, che
permetteva di contenere soltanto cose per Unreal, invece poi l'hanno reso più
ampio''. Ci sono quindi pacchetti con modelli 3D ``proprio grezzi'', da
sistemare. Se si vogliono pacchetti già pronti, con i materiali e tutto
``setappato bene'', conviene:
\begin{itemize}
  \item \textbf{filtrare per Unreal Engine} come formato;
  \item filtrare per \textbf{gratuito} (o ordinare per prezzo crescente:
        ``è tutto quanto gratis''), ``tanto c'è tanta roba gratis''.
\end{itemize}

\paragraph{I pacchetti gratuiti a tempo.} Ogni \textbf{quindici giorni} Fab
rende gratuiti \textbf{tre pacchetti}: ``se li riscattate, poi sono gratis per
sempre''. Si trovano in homepage; quella metà mese ci sono ``un robot
gigantesco, un tunnel e un launchpad''.

\paragraph{Non esagerare con il peso.} ``Non prendete pacchetti che pesano mille
gigabyte'': rallentano lo sviluppo, il download e tutto il resto. E, soprattutto:
``l'importante comunque per il task è \textbf{il codice}, non
l'\textbf{estetica}''. Il docente lascia un quarto d'ora per cercare i
pacchetti, ``non per farci troppo tempo: un pacchetto qualsiasi va bene''.

\subsection{Le versioni: avanti sì, indietro mai}

``Cosa super importante con i pacchetti e tutti quanti gli asset in generale.''
Gli asset di Unreal si possono sempre portare \textbf{da versioni precedenti a
versioni successive}, ma \textbf{mai il contrario}. Se si trovano pacchetti, o
si creano progetti, in una versione successiva, non si può più tornare indietro,
``perché si corrompe tutto'': ``gli asset scompaiono, fate proprio incolla e non
esistono più''.

\paragraph{Lo spazio su disco.} I pacchetti scaricati restano anche nella
\textbf{cache del Vault} sul computer: se ne può cancellare quella copia per
recuperare spazio, ``tuttavia lo stanno cambiando proprio adesso''.

\subsection{Un piccolo trucco per le mappe grandi}

Mentre il pacchetto si scarica (``che internet rapido che abbiamo''), il docente
segnala un trucco per aprire più in fretta le \textbf{mappe grandi}: invece del
doppio clic diretto, fare prima \textbf{clic destro} sulla mappa, così l'editor
comincia a caricarne il contenuto mentre si può ancora usare l'engine, e poi il
doppio clic la apre più rapidamente: ``un piccolo escamotage''.\footnote{La
trascrizione di questo passaggio è frammentaria; si riporta il senso del
suggerimento così come è ricostruibile.}

\paragraph{Nell'attesa.} Mentre gli studenti cercano i pacchetti, il docente
gira tra i banchi. Uno studente non trova i contenuti nel progetto: aveva creato
il progetto con la propria versione e Fab non li cercava lì; il consiglio è
aggiungerli dall'interno dell'editor, con \emph{Window $\rightarrow$ Fab}. A un
altro che chiede come far saltare il personaggio da codice, il docente indica il
nodo \textbf{Jump} (se ne riparla nella sezione~\ref{gd12:sec:trampolino}). E a
tutti ripete: ``mi raccomando, lo vedremo sul codice, non sull'estetica''.

% ---------------------------------------------------------------------
\section{La timeline}
\label{gd12:sec:timeline}

\subsection{Il nodo più complesso, e uno dei più utili}

Arriva il nodo che qualcuno ha già provato da solo: la \textbf{timeline}. ``In
realtà è molto comodo, perché permette di fare un sacco di cose.'' È il nodo
\textbf{con più pin} di Unreal, il più complesso, ``perché in realtà è tipo un
nodo componente''; però ``è super utile, si usa veramente tanto, ed è una delle
funzionalità veramente forti, secondo me, dei Blueprint, che gli altri linguaggi
di programmazione non hanno''. E una rassicurazione, che il docente ripeterà
poi a una studentessa in difficoltà: ``la timeline è tipo il nodo più complesso
che c'è: se capite bene questo, tutti gli altri sono più facili''.

\subsection{La piattaforma mobile}

Il docente crea un nuovo \textbf{Actor}, \textbf{MovingPlatform}, e ci mette una
piattaforma: un \textbf{Plane}.

\paragraph{Che cosa fa una timeline.} Permette di \textbf{spostare un oggetto, o
modificare qualsiasi parametro, nel tempo}, ``senza troppi problemi e con tante
funzionalità veramente interessanti''. Si crea cercando \emph{timeline} nel
grafo: \textbf{Add Timeline}. Non si può mettere dappertutto (``ci sono delle
limitazioni''\footnote{Per esempio, una timeline non si può aggiungere
all'interno di una \emph{funzione}: vive nell'Event Graph di un Blueprint, come
un componente.}), ma nell'Event Graph di un Actor sì. Si può, ``anzi si deve'',
\textbf{rinominare}: il docente la chiama \textbf{Movement}.

\paragraph{I pin.} La timeline fa \textbf{progredire dei valori} nel tempo, sia
avanti che indietro. A sinistra ha diversi ingressi di esecuzione:
\textbf{Play} (parte da dove si trova), \textbf{Play from Start} (riparte
dall'inizio), \textbf{Stop}, \textbf{Reverse} (torna indietro da dove si
trova), \textbf{Reverse from End}, \textbf{Set New Time}. A destra, le uscite:
\textbf{Update}, \textbf{Finished}, \textbf{Direction} e, per ogni traccia
creata, un pin con il suo valore. Collegata al \textbf{Begin Play} su Play, la
timeline partirà appena l'oggetto viene creato; ``tuttavia non sappiamo ancora
quali dati ci sono qua dentro''.

\subsection{L'editor della timeline}

Con il \textbf{doppio clic} sul nodo si apre un'altra schermata, quella in cui
si impostano \textbf{le variazioni nel tempo} di una variabile.
\begin{enumerate}
  \item Si aggiunge una \textbf{traccia} con il pulsante \textbf{+}. Quella che
        si usa ``il 99 per cento delle volte'' è la \textbf{Float Track}: un
        valore float (``dalla volta scorsa: variabile con la virgola'') che
        cambia nel tempo. Le si dà un nome.
  \item Con il \textbf{clic destro} $\rightarrow$ \textbf{Add Key} si aggiunge
        un punto, cioè un valore in un istante. Con un solo punto, il valore
        (per esempio 0,27) ``rimarrà a 0,27 per tutta la durata della
        timeline''.
  \item Con due punti si crea una \textbf{curva}, in questo caso una retta; con
        più punti si possono creare variazioni interessanti.
  \item Selezionando i punti si può scegliere l'interpolazione \textbf{Auto},
        così che il passaggio sia \textbf{curvo}, ``non soltanto dei valori
        netti''.
\end{enumerate}

\paragraph{Una piattaforma che sale e scende.} Per la piattaforma: al tempo
\textbf{0} il valore è \textbf{0} (``parte da per terra''); dopo \textbf{2,5
secondi} vale \textbf{200}, l'altezza che si vuole raggiungere; al secondo
\textbf{5}, l'ultimo istante della timeline, \textbf{torna a 0}. ``Adesso è
solamente un valore che fa su e giù, però dovrebbe farlo poi di conseguenza la
piattaforma.''

\begin{figure}[htbp]
\centering
\begin{tikzpicture}[x=2cm, y=0.016cm, font=\small]
  \draw[gray!30] (0,0) grid[xstep=0.5, ystep=50] (5,220);
  \draw[-{Stealth}] (0,0) -- (5.35,0) node[right] {tempo (s)};
  \draw[-{Stealth}] (0,0) -- (0,245) node[above] {valore};
  \foreach \t in {0,1,2,2.5,3,4,5} \node[below, font=\footnotesize] at (\t,0) {\t};
  \foreach \v in {0,100,200} \node[left, font=\footnotesize] at (0,\v) {\v};
  \draw[very thick, orange!80!black] (0,0) -- (2.5,200) -- (5,0);
  \draw[thick, dashed, blue!60!black] plot[smooth, domain=0:5, samples=60]
       (\x,{100-100*cos(deg(\x*pi/2.5))});
  \foreach \p in {(0,0),(2.5,200),(5,0)} \fill[black] \p circle (2.2pt);
  \node[font=\footnotesize, align=left, anchor=west] at (5.25,120)
       {\textcolor{orange!80!black}{\rule[2pt]{14pt}{1.2pt}} interpolazione lineare\\
        \textcolor{blue!60!black}{- - -} interpolazione \emph{Auto}};
  \draw[densely dotted] (5,0) -- (5,230) node[above, font=\footnotesize] {Length = 5};
\end{tikzpicture}
\caption{La Float Track della piattaforma: tre chiavi, $(0, 0)$, $(2{,}5, 200)$
e $(5, 0)$. Con l'interpolazione lineare il moto ha degli spigoli; con
\emph{Auto} la curva rallenta in cima e in fondo, e il movimento è più
morbido.}
\label{gd12:fig:curva}
\end{figure}

\paragraph{La lunghezza.} La durata della timeline è data dal campo
\textbf{Length}: ``potete dirgli: questa timeline dura 5 secondi, 20 secondi,
80 secondi, quello che volete''.

\paragraph{Un pin per traccia.} Aggiungendo la traccia, sul nodo è comparso un
nuovo pin con il suo valore: ``cosa molto comoda, perché in realtà ne possiamo
creare più di una''. Se serve che la piattaforma \textbf{si alzi e ruoti}, si
fanno due tracce diverse, e la timeline avrà due uscite.

\subsection{Collegare la curva alla piattaforma}

Dal componente Plane si tira fuori un \textbf{Set Relative Location}
(``poi vi spiego come mai''), si divide il vettore e si collega il valore della
traccia sulla $z$, ``dal basso verso l'alto''. L'esecuzione parte
dall'\textbf{Update}, che viene chiamato \textbf{a ogni frame} mentre la
timeline sta andando, ``quindi praticamente costantemente''.

\begin{figure}[htbp]
\centering
\begin{bpgraph}
  \bpnode{bp}{0,0}{bpHEvent}{Event BeginPlay}{2.6}{}{exec/}
  \bpnode{tl}{3.6,0}{bpHFlow}{Movement}{3.6}%
    {exec/Play, exec/Play from Start, exec/Stop, exec/Reverse, exec/Reverse from End, exec/Set New Time, float/New Time}%
    {exec/Update, exec/Finished, enum/Direction, float/Height}
  \bpget{pl}{4.2,-3.9}{obj}{Plane}{1.6}
  \bpnode[2]{sl}{10.6,0}{bpHFunc}{Set Relative Location\\\bpsub{Target is Scene Component}}{4.2}%
    {exec/, obj/Target, float/New Location X \bpval{0}, float/New Location Y \bpval{0}, float/New Location Z, bool/Sweep}{exec/}
  \bpwire{exec}{bp-o1}{tl-i1}
  \bpwire{exec}{tl-o1}{sl-i1}
  \bpwire{obj}{pl-o1}{sl-i2}
  \bpwire{float}{tl-o4}{sl-i5}
  \bpnote{7.6,-2.3}{Update: a ogni frame\\finché la timeline va}
\end{bpgraph}
\caption{La piattaforma mobile: al Begin Play la timeline \emph{Movement}
parte; a ogni Update il valore della traccia (qui chiamata \emph{Height})
diventa la $z$ relativa del Plane. Il pin della traccia compare da solo,
appena la traccia viene creata.}
\label{gd12:fig:piattaforma}
\end{figure}

Messa la piattaforma nella scena e premuto Play, ``vedete che si muove in su''.
(Un piano visto da sotto non si vede: il docente sta per un attimo ``sotto'' la
piattaforma e non la vede più, ma il movimento c'è.)

\paragraph{Il loop.} Nell'editor della timeline c'è l'opzione \textbf{Loop}:
attivata, la timeline riparte da sola all'infinito: ``fa su, giù, poi posso
salire, su, giù''. ``Abbastanza facile.''

\paragraph{Finished.} L'uscita \textbf{Finished} viene chiamata quando la
timeline arriva in fondo, ``nel caso in cui non la mettiate in loop,
chiaramente, altrimenti non finisce''.

\subsection{La porta: Play e Reverse}

``Una delle cose comode è che potete fare che avanza e poi torna indietro.''
Esempio: una \textbf{porta} che si apre quando ci si avvicina e si richiude
quando ci si allontana. ``Lo faccio anche molto rapidamente, tanto è super
facile.''
\begin{enumerate}
  \item Un \textbf{cubo} (``questo fantastico cubo sarà la nostra porta''),
        con un \textbf{Set Relative Rotation} (``relative, sempre''), e il
        valore della traccia collegato allo \textbf{yaw}, la rotazione attorno
        all'asse verticale: ``facciamo una porta praticamente girevole''.
  \item Si toglie il loop (``non è che si deve chiudere in automatico'') e si
        mette una \textbf{Box Collision} che fa da trigger, ingrandendola
        (``è un po' piccolina'').
  \item Sul \textbf{Begin Overlap} della box: \textbf{Play}. Sull'\textbf{End
        Overlap}: \textbf{Reverse}.
  \item La curva va da \textbf{0 a 90} gradi, ``non a 200'': la rotazione
        parte da zero, arriva a 90 (la porta si apre), poi con il Reverse torna
        indietro (la porta si chiude).
\end{enumerate}

\begin{figure}[htbp]
\centering
\begin{bpgraph}
  \bpnode[2]{bo}{0,0}{bpHEvent}{On Component Begin Overlap\\\bpsub{(Box)}}{4.0}{}{exec/, obj/Other Actor}
  \bpnode[2]{eo}{0,-2.2}{bpHEvent}{On Component End Overlap\\\bpsub{(Box)}}{4.0}{}{exec/, obj/Other Actor}
  \bpnode{tl}{5.4,0}{bpHFlow}{Door}{3.4}%
    {exec/Play, exec/Play from Start, exec/Stop, exec/Reverse, exec/Reverse from End}%
    {exec/Update, exec/Finished, enum/Direction, float/Angle}
  \bpnode[2]{sr}{10.4,0}{bpHFunc}{Set Relative Rotation\\\bpsub{Target is Scene Component}}{4.0}%
    {exec/, obj/Target, float/Roll \bpval{0}, float/Pitch \bpval{0}, float/Yaw}{exec/}
  \bpget{cu}{7.0,-3.3}{obj}{Cube}{1.5}
  \bpwire{exec}{bo-o1}{tl-i1}
  \bpwire{exec}{eo-o1}{tl-i4}
  \bpwire{exec}{tl-o1}{sr-i1}
  \bpwire{obj}{cu-o1}{sr-i2}
  \bpwire{float}{tl-o4}{sr-i5}
\end{bpgraph}
\caption{La porta girevole: entrare nella box fa avanzare la timeline (da 0 a
90 gradi di yaw), uscirne la fa tornare indietro dal punto in cui si trova. Con
Reverse la porta non ``scatta'': se si esce a metà apertura, si richiude da
lì.}
\label{gd12:fig:porta}
\end{figure}

\paragraph{Il pivot.} Il primo risultato ``è un pochettino strano'': la porta
ruota attorno al proprio \textbf{centro}, non attorno ai cardini, ``perché devo
ruotare il pivot, non ogni cosa''. Il docente sistema il pivot, e la porta si
apre e si chiude come deve: ``così si apre, così si chiude. Questa è l'idea''.
(Sul pivot si veda il capitolo~\ref{gd11:cap:blockout-luci}; un modo più pulito
per ottenere lo stesso effetto, con le \emph{scene component}, è nella
sezione~\ref{gd12:sec:doppia-porta}.)

\paragraph{A che cosa serve, in generale.} ``In realtà si usa per tantissime
cose la timeline'': per tutto ciò che \textbf{varia nel tempo}, ``la timeline
torna comoda''. E c'è un motivo in più per capirla bene: lo stesso strumento,
``un pochettino più complesso'', è quello delle \textbf{cinematiche}
(\emph{Sequencer}), ``perché è praticamente lo stesso tool''.

\paragraph{Approfondimento: un esempio dal libro.} Romero e Sewell introducono
la timeline con un esempio simile per spirito: lo \textbf{zoom} di un'arma,
in cui il \emph{Field of View} della telecamera passa da 90 a 45 gradi in 0,3
secondi. Il libro mostra due dettagli pratici che tornano utili anche qui: nel
riquadro in alto a sinistra del grafico si possono \textbf{scrivere a mano
tempo e valore} di un punto selezionato, per essere precisi, e due piccoli
pulsanti \textbf{adattano lo zoom} del grafico, così che i punti non escano
dalla vista. E fa notare il vantaggio della timeline rispetto a un SET secco:
il passaggio da un valore all'altro è \emph{graduale}, non uno scatto.

% ---------------------------------------------------------------------
\section{Relative e world}
\label{gd12:sec:relative}

``Cosa che non ho approfondito.'' Cercando \emph{Set Location} escono due
versioni: \textbf{Set Relative Location} e \textbf{Set World Location} (lo
stesso vale per rotazione e scala). Che differenza c'è?
\begin{itemize}
  \item \textbf{Relative}: relativo al componente rispetto all'\emph{actor} in
        cui si trova, ``quindi proprio queste coordinate qui di lato'', quelle
        del pannello dei dettagli del componente.
  \item \textbf{World}: ``sposta magari in questo punto del mondo,
        indipendentemente da dove sia l'oggetto''.
\end{itemize}
Per una piattaforma o una porta che si muovono \emph{rispetto a dove le si
mette}, quasi sempre si vuole \textbf{relative}: si può piazzare la stessa
piattaforma in dieci punti del livello, e ognuna salirà di 200 a partire da dove
si trova. Con world, tutte andrebbero allo stesso punto del mondo. ``Quindi fate
attenzione a questo, semplicemente.'' (Il collegamento con le coordinate locali
e globali del capitolo~\ref{gd11:cap:blockout-luci} è diretto.)

\paragraph{Location non è transform.} A una domanda il docente precisa che il
vettore di un Set Location contiene \textbf{solo la posizione}, non la
rotazione: ``non è una transform. Se fosse transform, includerebbe anche la
rotazione e la scala''. Per cambiare insieme posizione e rotazione servono due
nodi diversi, oppure un nodo che lavora su tutta la transform.

% ---------------------------------------------------------------------
\section{Il giro tra i banchi, prima parte}
\label{gd12:sec:giro1}

``Vi lascio un attimino fare, perché chiaramente ho fatto molte cose e potrebbero
non essere tutte chiare.'' Il docente lascia il proprio codice a schermo
(``così, se volete prenderne ispirazione, potete farlo tranquillamente'') e
passa a guardare gli schermi.

\subsection{La porta con le box di overlap}

Una studentessa sta costruendo la porta. Il docente la guida passo per passo.
\begin{itemize}
  \item \textbf{Gli eventi della box.} Selezionata la box, si scorre il
        pannello dei dettagli ``giù in fondo in fondo'': lì ci sono i
        pulsanti \textbf{+} degli eventi. \emph{On Component Begin Overlap}
        ``si triggera quando entra qualcosa'', \emph{On Component End Overlap}
        (il terzo) quando esce. ``Tendenzialmente il player, ma potrebbe
        entrare anche altro.''
  \item \textbf{La box è più piccola del cubo.} ``Mi sa che è più piccola del
        cubo stesso, quindi è tipo dentro'': nella viewport del Blueprint va
        resa \textbf{più grande della mesh}.
  \item \textbf{La box è figlia del cubo.} Così ``si muoverà con il cubo'':
        se non lo si vuole, non deve essere \emph{child}. Nella gerarchia dei
        componenti (``rimani sulla parte component'') la si trascina
        sulla radice, così i due sono ``slegati''. ``Avevo intenzione di farlo,
        ma mi sembrava fosse sbagliato.'' ``No, era la cosa giusta.''
  \item \textbf{Il nome della timeline.} La sua si chiama ancora
        \emph{Timeline}: è il nome di default, ``poi la devi rinominare''.
  \item \textbf{Il flusso.} Collegata al Begin Play, ``entra nella timeline, e
        da quando entra andrà tutti quanti i frame finché la timeline non
        finisce''.
  \item \textbf{La traccia.} Doppio clic, \emph{+ Float Track} (``perché
        vogliamo cambiare un valore float''), clic destro \emph{Add Key}: la
        linea che compare è ``come cambierà questo valore nel tempo''.
  \item \textbf{Vedere il valore.} Sull'Update, un \textbf{Print String} con
        il valore della traccia collegato al testo (``si converte''): avvicinandosi
        alla porta, ``vedi che c'è il valore che è cambiato rapidamente come gli
        hai detto tu prima''.
\end{itemize}
Un errore di compilazione sul \emph{ThirdPersonCharacter} (un nodo rimasto lì
``che non portava a niente'') blocca il primo Play: tolto quello, funziona.

\paragraph{Play non è Play from Start.} La timeline va solo \textbf{la prima
volta}: ``perché gli hai fatto Play e non Play from Start''. Play riprende da
dove la timeline si trova: arrivata in fondo, resta in fondo. Per farla
ripartire, o \textbf{Play from Start}, o la si fa tornare indietro con
\textbf{Reverse} e poi ripartire.

\subsection{Il teletrasporto a colpi di Delay}

Antonino\footnote{``Antonio'' nella trascrizione: con ogni probabilità
Antonino, presente nella revisione del capitolo~\ref{gd11:cap:blockout-luci}.}
condivide lo schermo: si è concentrato solo sui Blueprint, e ha creato un
\textbf{sistema di teletrasporto} che, quando lo si usa, teletrasporta il
personaggio e poi \textbf{lo ingrandisce}, con un \textbf{messaggio} in alto a
sinistra che avvisa che l'abilità si è disattivata, un \textbf{cooldown di 5
secondi}, e poi la riattivazione. ``Particolare, però può essere
interessante.''

Il codice è fatto \textbf{tutto con dei Delay}: ``mi piace il fatto che ti sia
impegnato molto; chiaramente ci sono soluzioni più eleganti, tuttavia vedo che
sei riuscito a raggiungere l'obiettivo che volevi. Quindi bravo.'' E con una
battuta: ``Tu hai creato una timeline con le tue forze: delay, delay, delay.''

\paragraph{Approfondimento: Delay contro timeline.} Una catena di Delay va
bene per \emph{aspettare} (il cooldown), ma è rigida per \emph{animare}: ogni
passo è uno scatto, e cambiare i tempi significa ritoccare ogni nodo. Una
timeline descrive l'intera evoluzione in un solo grafico, la interpola fra un
frame e l'altro, e si può riprodurre all'indietro o interrompere; per
l'ingrandimento graduale del personaggio sarebbe lo strumento naturale.

Mattia, invece, non riesce a condividere lo schermo (``mi ha freezato''): il
docente gli suggerisce di fare \emph{Zip Project} e mandarlo, così lo si guarda
dopo.

\subsection{La piattaforma che va avanti e indietro}

Un altro studente vuole una piattaforma che si muova \textbf{avanti e indietro},
a destra e sinistra, \emph{per sempre}. È l'occasione per vedere da vicino
alcuni errori tipici.
\begin{enumerate}
  \item \textbf{Prima di tutto, la parte fisica.} Si trascina nel grafo il
        componente che si vede (il cubo) e se ne tira fuori un \textbf{Set
        Relative Location}: ``per spostarla deve cambiare la posizione,
        immagino''.
  \item \textbf{Collegare l'Update.} ``Devi collegarla, se no non è
        collegata e chiaramente non può andare'': senza filo bianco il nodo non
        viene mai eseguito.
  \item \textbf{Dividere il vettore.} Collegato così, il nodo porta la
        posizione a $(0,0,0)$ a ogni frame, ``che non è quello che vuoi''. Con
        il clic destro sul pin \textbf{New Location} $\rightarrow$
        \textbf{Split Struct Pin} il vettore si divide nelle tre componenti, e
        il valore della traccia va sulla \textbf{$x$}, ``che è l'avanti e
        indietro''.
  \item \textbf{Attenzione all'origine.} Set Relative Location imposta
        \emph{proprio} quei valori: se il componente non parte da $(0,0,0)$,
        al primo frame ``ti si sposta di colpo''.
  \item \textbf{Valori sensati.} La piattaforma ``trema'' molto velocemente:
        i valori nella curva sono minuscoli, ``si muove di pochissimo, quindi
        potresti anche non accorgertene''. Il docente fa cancellare i punti e
        rifarli \textbf{scrivendo i valori a mano} nei campi \emph{Time} e
        \emph{Value}: $(0, 0)$, e poi a \textbf{5 secondi} il valore
        \textbf{200}. Ora si muove: ``qualcosa sta funzionando''.
  \item \textbf{Partire subito.} Lo studente vuole che vada all'infinito,
        ``anche senza attivarlo'': si cancella la parte con l'overlap, e la
        timeline si collega al \textbf{Begin Play}, ``così parte subito quando
        viene creato l'oggetto''.
  \item \textbf{Il loop.} Si attiva \textbf{Loop} nell'editor della
        timeline (``un pochettino più in alto, a sinistra'').
  \item \textbf{Lo scatto all'indietro.} Ora la piattaforma va avanti\dots{} e
        poi \emph{ricompare} all'inizio. ``È normale'': il valore parte da
        zero, sale fino a 200 e poi, ricominciando il loop, \emph{riparte da
        zero}. ``Al posto di fare una curva che è una linea, devi fare che va
        su e che poi torna giù'': un punto in mezzo, a \textbf{2,5 secondi}, al
        valore massimo, e quello a \textbf{5 secondi} riportato a
        \textbf{0}. ``Se lo metti a 0, torna indietro, perché la posizione
        torna a 0.''
\end{enumerate}
I due pulsanti in alto a sinistra del grafico (\emph{zoom to fit} orizzontale e
verticale) ridimensionano la vista, così che si vedano tutti i punti. ``Va
avanti, torna indietro, torna avanti. Direi che funziona.'' ``Ok, perfetto,
questa cosa mi servirà.'' Il consiglio finale: fare \textbf{diverse cose con la
timeline}, ``in modo da prenderci mano''.

\paragraph{Collisioni e overlap.} A uno studente che chiede se usare la mesh
stessa per gli overlap, il docente risponde che può essere comodo per forme
complesse (un arco, per esempio, controllando che non dia overlap strani), ma
che di solito per gli overlap si usa \textbf{una box}.

% ---------------------------------------------------------------------
\section{Muovere un oggetto con l'Event Tick}
\label{gd12:sec:tick}

Una studentessa chiede come far muovere qualcosa senza la timeline. ``Ci sono
tanti modi per far muovere qualcosa.'' Il primo è la timeline; il docente ne
mostra un altro, ``soprattutto per muovere le cose non da punto A a punto B, ma
in generale''.

\paragraph{Il modo base.} Un nuovo Actor, \textbf{MovingObject}, che deve
\textbf{salire sempre}. Sull'\textbf{Event Tick}, ``perché si deve aggiornare
tutti quanti i frame'': \textbf{Set Actor Location} con, come nuova
posizione, \textbf{Get Actor Location} \textbf{+} $(0, 0, 10)$. Così, a ogni
frame, l'oggetto sale di 10 centimetri, ``e quindi parte e vola''. (Senza
quel \emph{+}, ``sarebbe tutti quanti i frame sempre dov'è, e quindi non
servirebbe a nulla''.)

\paragraph{Funziona\dots{} quasi.} ``Chiaramente questa cosa qua funziona
\emph{circa}, perché va tutti quanti i frame, giusto, però ognuno ha un frame
rate diverso sul proprio PC.'' È il problema già visto nel
capitolo~\ref{gd7:cap:blueprint} (sezione~\ref{gd7:sec:delta}): il valore va
\textbf{moltiplicato per i Delta Seconds}. Con $10 \times \Delta t$ l'oggetto
sale di 10 centimetri \textbf{al secondo} (``magari anche un pochettino di più,
perché è un po' pochino''), su qualsiasi PC: ``vedete che piano piano si
solleva''.

\begin{figure}[htbp]
\centering
\begin{bpgraph}
  \bpnode{tk}{0,0}{bpHEvent}{Event Tick}{2.8}{}{exec/, float/Delta Seconds}
  \bpnode{mul}{3.6,-1.5}{compact}{$\times$}{1.8}{float/, float/\bpval{100}}{float/}
  \bpnode{mk}{6.2,-1.5}{bpHPure}{Make Vector}{2.4}{float/X \bpval{0}, float/Y \bpval{0}, float/Z}{vec/}
  \bpnode{ga}{6.0,-3.6}{bpHPure}{Get Actor Location}{3.2}{obj/Target \bpval{self}}{vec/Return Value}
  \bpnode{add}{9.8,-2.2}{compact}{$+$}{1.4}{vec/, vec/}{vec/}
  \bpnode[2]{sl}{11.8,0}{bpHFunc}{Set Actor Location\\\bpsub{Target is Actor}}{3.4}{exec/, obj/Target \bpval{self}, vec/New Location}{exec/}
  \bpwire{exec}{tk-o1}{sl-i1}
  \bpwire{float}{tk-o2}{mul-i1}
  \bpwire{float}{mul-o1}{mk-i3}
  \bpwire{vec}{mk-o1}{add-i1}
  \bpwire{vec}{ga-o1}{add-i2}
  \bpwire{vec}{add-o1}{sl-i3}
\end{bpgraph}
\caption{Il modo più semplice di muovere un oggetto: a ogni frame, la posizione
attuale più uno spostamento. Moltiplicare la velocità (qui 100 cm/s) per i
Delta Seconds rende il movimento uguale su tutti i computer.}
\label{gd12:fig:tick}
\end{figure}

\paragraph{La base sotto tutto il resto.} ``Questo è il modo più facile per far
muovere un oggetto nello spazio, il modo proprio base.'' La timeline e gli altri
strumenti sono ``tutte quante cose che noi mettiamo poi sopra per semplificare
un attimino la vita''. Ma se si vuole far muovere, per esempio, ``un
\textbf{proiettile} sempre dritto, si fa così''.

Durante la prova un'altra piattaforma, fatta con un piano, ``da sotto non si
vede'': ``metti un cubo e si risolve''. E sul PC del docente, molto più veloce,
la versione senza Delta Seconds andrebbe ``molto più veloce'': un'ulteriore
conferma del perché moltiplicare.

% ---------------------------------------------------------------------
\section{La piattaforma che si distrugge}
\label{gd12:sec:distrugge}

\subsection{Una porta che si apre da sola}

Prima, una correzione veloce a una studentessa che ha messo \textbf{due box}
davanti alla porta: la porta si apre anche quando nessuno si avvicina, perché
``gli stessi oggetti stanno entrando nella box quando si apre''. L'overlap
scatta per \emph{qualsiasi cosa} entri, anche per le ante della porta stessa.
La soluzione è quella del capitolo~\ref{gd7:cap:blueprint}
(sezione~\ref{gd7:sec:trigger}): controllare chi è entrato, con un
\textbf{Cast To ThirdPersonCharacter} dall'\emph{Other Actor}, oppure
confrontando l'Other Actor con il giocatore. ``L'Other Actor ti fornisce
l'informazione di che cosa sta toccando questa box'': così si è sicuri che sia
davvero il personaggio.

\subsection{Francesco: la piattaforma che crolla}

Francesco vuole una \textbf{piattaforma che si distrugge} quando ci si sale
sopra (accanto a un'altra, stabile). ``Bella come idea! Ci sta.'' Al primo Play,
però, la piattaforma \textbf{scompare subito}. ``Guarda il codice: Event Begin
Play, ecco perché'': il Destroy era collegato al Begin Play, ``quindi appena
parte si distrugge''.

\begin{enumerate}
  \item Il trigger giusto è l'\textbf{overlap}, ``con quella box che hai già
        lì'', non il Begin Play.
  \item Secondo errore: il nodo stava distruggendo \textbf{la box}, cioè il
        componente che controlla l'overlap. ``L'overlap è quello giusto, è il
        componente che distruggi che è sbagliato.'' Bisogna distruggere il
        \textbf{Plane}, oppure tutto l'\textbf{Actor} con \emph{Destroy
        Actor}, ``come stavi facendo prima''.
  \item \emph{Destroy Component} contro \emph{Destroy Actor}: il primo toglie
        un solo componente, il secondo ``ti distrugge tutto quanto''. Qui
        vanno bene entrambi, purché si distrugga la parte giusta.
\end{enumerate}
``Se di una cosa non sei sicuro, \textbf{provala} e vedi come va.'' Ora
funziona. E un'idea in più: mettere un piccolo \textbf{Delay} prima della
distruzione, ``se vuoi che si distrugga dopo'': si può restare un secondo sulla
piattaforma, poi crolla.

% ---------------------------------------------------------------------
\section{Ricaricare il livello}
\label{gd12:sec:ricarica}

``Ah, vi ricordo questa cosa \textbf{fondamentale per la riuscita del task}.''
Se si apre lo stesso livello in cui ci si trova, il livello \textbf{si riapre da
capo}, ``quindi vi refresha tutto''. C'è un modo ``super comodo'' per farlo
senza scrivere a mano il nome della mappa:
\begin{center}
\textbf{Get Current Level Name} $\rightarrow$ \textbf{Open Level (by Name)}
\end{center}
``In questo modo vi riapre lo stesso livello, quindi si resetta.'' Per esempio,
quando il giocatore muore, si fa questo e il livello si ricarica: ``molto
facile, però segnatevelo''. È esattamente ciò che serve per due requisiti del
task: la \textbf{trappola} che fa ricominciare il livello e il \textbf{timer}
che, se scade, lo fa ripartire.

\begin{figure}[htbp]
\centering
\begin{bpgraph}
  \bpnode[2]{bo}{0,0}{bpHEvent}{On Component Begin Overlap\\\bpsub{(Box)}}{4.0}{}{exec/, obj/Other Actor}
  \bpnode{ca}{5.0,0}{bpHCast}{Cast To ThirdPersonCharacter}{4.6}{exec/, obj/Object}%
    {exec/, exec/Cast Failed, obj/As BP Third Person Character}
  \bpnode[2]{gl}{6.0,-2.4}{bpHPure}{Get Current Level Name\\\bpsub{Remove Prefix String \bpval{true}}}{3.9}{}{str/Return Value}
  \bpnode{ol}{10.8,0}{bpHFunc}{Open Level (by Name)}{3.6}{exec/, name/Level Name, bool/Absolute, str/Options}{exec/}
  \bpwire{exec}{bo-o1}{ca-i1}
  \bpwire{obj}{bo-o2}{ca-i2}
  \bpwire{exec}{ca-o1}{ol-i1}
  \bpwire{str}{gl-o1}{ol-i2}
\end{bpgraph}
\caption{La trappola che fa ricominciare il livello: se a entrare nella box è
il personaggio, si legge il nome del livello corrente e lo si riapre. La
stringa viene convertita automaticamente nel tipo \emph{Name} del pin.}
\label{gd12:fig:ricarica}
\end{figure}

\paragraph{Il livello che si riapre all'infinito.} Più tardi, uno studente
inciampa proprio qui: il gioco si riapre di continuo. ``Stai riaprendo lo stesso
livello all'infinito.'' L'Open Level era collegato a un overlap che scattava
\emph{per qualsiasi cosa}, anche per oggetti già dentro la box all'avvio:
appena il livello si apriva, l'overlap lo riapriva. La cura è la stessa di prima:
un \textbf{cast al personaggio}, per controllare che a toccare sia davvero il
giocatore, e solo allora ricaricare.

\paragraph{Approfondimento: che cosa si perde ricaricando.} Riaprire il livello
lo ricostruisce da zero: ogni Actor torna allo stato iniziale, e con lui
\emph{tutte le variabili} che contiene (monete raccolte, tempo trascorso). È
esattamente l'effetto voluto per una trappola, ma vuol dire anche che un dato
che deve sopravvivere al reset (per esempio un contatore delle morti, o il
punteggio totale fra i tre livelli) non può stare su un Actor del livello: va
messo nella \emph{Game Instance}, come anticipato nel
capitolo~\ref{gd7:cap:blueprint} (sezione~\ref{gd7:sec:vita}).

% ---------------------------------------------------------------------
\section{La doppia porta}
\label{gd12:sec:doppia-porta}

\subsection{Il problema}

Una studentessa vuole una porta \textbf{a due ante} che si apra \textbf{nella
direzione giusta}: verso l'interno se si arriva da davanti, verso l'esterno se
si arriva da dietro. ``Interessante, quello. Si può fare.'' Il consiglio:
\textbf{due box}, una davanti e una dietro; in base a quale delle due viene
toccata, le ante si aprono in un verso o nell'altro.

I primi tentativi incontrano i problemi classici. La porta si apre e richiude
male se si passa da una box all'altra prima che l'animazione sia finita
(``quando triggeri quella davanti, poi vai nell'altra box e ti triggeri quella
dopo, prima che si sia chiusa l'altra''). Il docente suggerisce una variabile
\textbf{booleana} che dica se le porte sono chiuse o se l'animazione è finita:
``di base è vera, quindi le porte sono chiuse''. Quando ci si avvicina la si
mette a falso, e la si rimette a vero quando la porta si è richiusa. Attenzione
all'ordine: ``tu la stai settando come vera e poi la controlli subito, fa poco
senso. Prima controlli, e poi in caso la setti''. Con un \textbf{Sequence} si
possono ordinare le operazioni, e con un \textbf{Delay} si aspetta che
l'animazione finisca. Il docente cita anche \textbf{Combine Rotators}, che somma
due rotazioni (``come se fosse una somma, ma per le rotazioni'').

Ma con due box e due timeline separate il codice ``va un po' fuori controllo''.
``Facciamo così, aspetta un attimo, vediamo come si potrebbe rilavorare'': il
docente la costruisce da capo, in diretta.

\subsection{La versione del docente}

Il nuovo Blueprint, \emph{DoubleDoor}, contiene \textbf{due box} (davanti e
dietro) e \textbf{due cubi} (le ante, \emph{Cube1} e \emph{Cube2}). Una sola
\textbf{timeline} porta un valore da \textbf{0 a 90 in 2,5 secondi}. Due
variabili booleane:
\begin{itemize}
  \item una che dice se \textbf{le animazioni sono finite} (nella lezione
        \emph{is animation}), con valore di default \textbf{vero}: finché è
        falsa, non si può rientrare;
  \item \textbf{IsInFront}, che ricorda da quale lato è arrivato il
        personaggio.
\end{itemize}

\begin{figure}[htbp]
\centering
\begin{tikzpicture}[
  font=\small,
  gd12box/.style={draw, thick, rounded corners=3pt, align=center,
                  minimum height=0.8cm, inner sep=3pt, text width=3.0cm},
  gd12ev/.style={gd12box, fill=red!12},
  gd12dec/.style={draw, thick, diamond, aspect=2.2, align=center,
                  inner sep=1pt, fill=gray!10, font=\footnotesize},
  gd12tl/.style={gd12box, fill=blue!8},
  gd12ar/.style={-{Stealth[length=2.3mm]}, thick}
]
  \node[gd12ev] (front) at (0,0) {\textbf{Begin Overlap}\\box davanti};
  \node[gd12ev] (back) at (6,0) {\textbf{Begin Overlap}\\box dietro};
  \node[gd12dec] (c1) at (0,-1.6) {animazioni\\finite?};
  \node[gd12dec] (c2) at (6,-1.6) {animazioni\\finite?};
  \node[gd12box] (s1) at (0,-3.2) {IsInFront = \textbf{vero}\\finite = falso};
  \node[gd12box] (s2) at (6,-3.2) {IsInFront = \textbf{falso}\\finite = falso};
  \node[gd12tl] (tl) at (3,-4.8) {\textbf{Timeline}\\Play from Start\\(0 $\rightarrow$ 90 in 2,5 s)};
  \node[gd12dec] (dir) at (3,-6.6) {Direction\\= Forward?};
  \node[gd12box] (rev) at (-1.3,-6.6) {\textbf{Delay} 2 s\\poi \textbf{Reverse}};
  \node[gd12box] (fin) at (3,-8.4) {finite = \textbf{vero}\\(si può rientrare)};
  \node[gd12box, text width=3.9cm] (up) at (7.2,-6.6) {\textbf{Update}:\\yaw anta 1 $= v \cdot s$\\yaw anta 2 $= -v \cdot s$\\{\footnotesize $s=+1$ se IsInFront, altrimenti $-1$}};
  \draw[gd12ar] (front) -- (c1);
  \draw[gd12ar] (back) -- (c2);
  \draw[gd12ar] (c1) -- node[right, font=\footnotesize] {sì} (s1);
  \draw[gd12ar] (c2) -- node[right, font=\footnotesize] {sì} (s2);
  \draw[gd12ar] (s1) |- ([yshift=0.2cm]tl.west);
  \draw[gd12ar] (s2) |- ([yshift=0.2cm]tl.east);
  \draw[gd12ar, dashed] ([yshift=-0.25cm]tl.east) -| node[above, font=\footnotesize, pos=0.3] {ogni frame} (up.north);
  \draw[gd12ar] (tl) -- node[right, font=\footnotesize] {Finished} (dir);
  \draw[gd12ar] (dir) -- node[above, font=\footnotesize] {sì: aperta} (rev);
  \draw[gd12ar] (dir) -- node[right, font=\footnotesize] {no: chiusa} (fin);
  \draw[gd12ar] (rev.north) |- ([yshift=-0.25cm]tl.west);
\end{tikzpicture}
\caption{La logica della doppia porta costruita a lezione. Una sola timeline
guida entrambe le ante; il segno $s$, scelto in base al lato da cui si arriva,
decide se si aprono verso l'interno o verso l'esterno, e la seconda anta usa il
segno opposto alla prima. Quando la timeline finisce in avanti, dopo due
secondi torna indietro (la porta si richiude); quando finisce all'indietro, la
porta è di nuovo utilizzabile.}
\label{gd12:fig:doppia-porta}
\end{figure}

\paragraph{La spiegazione del docente.} ``Ve lo spiego un attimino, nel caso in
cui qualcuno volesse replicare questa mostruosità.''
\begin{enumerate}
  \item Prima si controlla, con la variabile booleana, che \textbf{non ci siano
        animazioni in corso}; di base è vera, quindi le animazioni sono finite
        e si può entrare.
  \item Se si entra nella box \textbf{davanti}, si mette \textbf{IsInFront} a
        vero (nella box dietro, a falso), e si mette l'altra variabile a
        \textbf{falso}, ``in modo che non possa più rientrare un'altra volta''.
  \item Si fa partire la \textbf{timeline}. Quando finisce, l'uscita
        \textbf{Direction} dice se stava andando \textbf{in avanti} o
        \textbf{in reverse}.
  \item Se stava andando in avanti, la porta si è aperta: si aspettano
        \textbf{due secondi} e poi la si fa chiudere, cioè ``riparte tutta
        quanta la timeline, ma in reverse''.
  \item Quando finisce il reverse, l'animazione è davvero finita: la variabile
        torna \textbf{vera} e si può rientrare.
  \item Durante l'Update, in base a IsInFront, il valore viene usato così com'è
        oppure \textbf{moltiplicato per $-1$}; e la seconda anta usa il segno
        opposto della prima, ``quindi una va di là e l'altra va di qua''.
\end{enumerate}
Al primo Play le ante si aprono al contrario (``vabbè, non è grave, si può
risolvere in fretta'') e una delle due non era collegata (``ho collegato sempre
lo stesso cubo''); sistemati i collegamenti e i segni, ``se esco si chiude, se
vado di qua mi aspetto che si apra al contrario, e infatti funziona''.

``Non bellissimo, non semplicissimo.'' E soprattutto: ``è chiaro che non mi
aspetto che facciate codice di questo tipo alla seconda lezione''. Chi vuole può
copiarlo; ``se avete dubbi, chiedetemi pure''.

\paragraph{Approfondimento: perché una sola timeline.} La versione del docente
risolve il caos dei tentativi precedenti con due scelte. La prima è avere
\textbf{un'unica fonte di verità} per lo stato della porta (una timeline,
un booleano che fa da \emph{lucchetto}) invece di due timeline indipendenti che
possono contraddirsi. La seconda è separare \textbf{che cosa} si anima (l'angolo
da 0 a 90) da \textbf{come} lo si applica (il segno, deciso dal lato): la curva
resta una sola, e la direzione è un parametro. È lo stesso principio che, più
avanti, porterà alle \emph{state machine}.

\subsection{Le cerniere: le scene component}

``Consiglio personale.'' Per porte che sembrino porte vere, il docente aggiunge
\textbf{due Scene Component}\footnote{Nella trascrizione ``due sin'': sono le
\emph{Scene Component}, componenti vuoti che hanno solo una transform.} e ci
attacca sotto i cubi, uno per scena.
\begin{itemize}
  \item Le scene sono ``praticamente i \textbf{punti di ancoraggio}'':
        ``semplicemente sono degli elementi vuoti''.
  \item Si posizionano le scene \textbf{sull'angolo} dove dovrebbero esserci i
        cardini, e si ruotano \textbf{le scene al posto dei cubi}.
  \item I cubi, essendo figli, ruotano attorno a quel punto: la porta ``si apre
        proprio sui bordi delle porte, non al centro''.
\end{itemize}
``Funziona praticamente come se fosse il pivot''; è il modo più pulito di
ottenere ciò che, con la prima porta, richiedeva di spostare il pivot della mesh.
``Tuttavia qua stiamo andando un pochettino troppo avanti.''

% ---------------------------------------------------------------------
\section{Engine Scalability: se il PC non regge}
\label{gd12:sec:scalability}

``Informazione di servizio.'' Vedendo che diversi PC degli studenti faticano, il
docente mostra come abbassare la qualità della grafica \emph{dell'editor}:
\emph{Settings $\rightarrow$ \textbf{Engine Scalability Settings}} e poi
\textbf{Low}. ``È come se metteste le impostazioni basse sulla grafica'' di un
gioco: la viewport è più brutta, ma tutto è più fluido. (``L'hanno spostato 25
volte, il tasto'': in versioni diverse si trova in punti leggermente diversi
della barra.)

% ---------------------------------------------------------------------
\section{La lava}
\label{gd12:sec:lava}

Francesco vuole importare dei modelli 3D per creare un \textbf{terreno su cui
si cade}: ``c'è della lava e, appena tocchi la lava, muori''. Il consiglio
generale: \textbf{prima il codice}, poi lo si mette sulla mappa ``appena puoi''.
(È lo stesso principio del blockout del capitolo~\ref{gd11:cap:blockout-luci}:
prima le funzionalità, poi l'aspetto.)

Il codice non compila. ``Ti manca veramente poco: è solo un nodo che non è
collegato.'' Il \textbf{Cast To ThirdPersonCharacter} ha il pin
\textbf{Object} vuoto, ``quindi non sa che cosa deve controllare che sia un
ThirdPersonCharacter''. Va collegato l'\textbf{Other Actor} dell'overlap,
``cioè l'attore che sta collidendo con questo cubo''. Errore successivo, e
classico: Francesco stava controllando se \emph{la box} fosse il personaggio;
va controllato l'Other Actor che arriva dall'evento. Poi lo si uccide
distruggendolo: ``direi che ha funzionato''. (Per il task, naturalmente, al posto
della distruzione si ricarica il livello, sezione~\ref{gd12:sec:ricarica}.)

``Finché non hai compilato, non puoi neanche lavorare'': un Blueprint con errori
di compilazione non gira, e l'editor lo segnala sul pulsante \emph{Compile} e
nel pannello dei risultati.

% ---------------------------------------------------------------------
\section{Il timer: sommare i Delta Seconds}
\label{gd12:sec:timer}

Uno studente chiede come costruire il \textbf{timer} del task: ``c'è il delta
time, ma è solo il delta time\dots{}''. Il docente parte dalla definizione:
i \textbf{Delta Seconds} sono già \emph{i secondi di un frame}, cioè quanto tempo
è passato dal frame precedente. Quindi basta \textbf{sommarli}.
\begin{enumerate}
  \item Si crea una variabile float, \emph{Value}.
  \item Sull'Event Tick: \emph{Value} $=$ \emph{Value} $+$ \textbf{Delta
        Seconds}.
  \item Con un Print String si vede che il valore ``aumenta di un secondo per
        ogni secondo''.
\end{enumerate}
Attenzione a non confondersi: se a ogni frame si somma \textbf{1}, invece, si
contano \textbf{i frame}, non i secondi.

\begin{figure}[htbp]
\centering
\begin{bpgraph}
  \bpnode{tk}{0,0}{bpHEvent}{Event Tick}{2.6}{}{exec/, float/Delta Seconds}
  \bpget{gv}{0.5,-2.4}{float}{Time}{1.4}
  \bpnode{add}{2.9,-1.8}{compact}{$+$}{1.3}{float/, float/}{float/}
  \bpnode[1]{set}{4.4,0}{bpHVar}{SET}{1.8}{exec/, float/Time}{exec/, float/}
  \bpnode{ge}{6.5,-1.6}{compact}{$\geq$}{1.6}{float/, float/\bpval{60}}{bool/}
  \bpnode{br}{8.3,0}{bpHFlow}{Branch}{2.1}{exec/, bool/Condition}{exec/True, exec/False}
  \bpnode{ol}{10.9,0}{bpHFunc}{Open Level (by Name)}{3.3}{exec/, name/Level Name}{exec/}
  \bpnode{gl}{6.2,-3.0}{bpHPure}{Get Current Level Name}{3.6}{}{str/Return Value}
  \bpwire{exec}{tk-o1}{set-i1}
  \bpwire{float}{gv-o1}{add-i1}
  \bpwire{float}{tk-o2}{add-i2}
  \bpwire{float}{add-o1}{set-i2}
  \bpwire{exec}{set-o1}{br-i1}
  \bpwire{float}{set-o2}{ge-i1}
  \bpwire{bool}{ge-o1}{br-i2}
  \bpwire{exec}{br-o1}{ol-i1}
  \bpwire{str}{gl-o1}{ol-i2}
\end{bpgraph}
\caption{Un timer di livello: a ogni frame si aggiungono i Delta Seconds alla
variabile \emph{Time}, che quindi conta i secondi reali; quando supera il limite
(qui 60 secondi), il livello riparte. La lezione mostra la somma; il confronto
e il ricaricamento sono il completamento naturale per il requisito del task.}
\label{gd12:fig:timer}
\end{figure}

\paragraph{Tanti modi.} ``Se non vi viene, potete anche calcolare il timer con
una timeline o qualsiasi roba che vi venga in mente. Non è un problema. Ci sono
tanti modi: per ogni cosa ci sono un sacco di soluzioni.''

\paragraph{Approfondimento: il timer di sistema.} Oltre alla somma sull'Event
Tick, Unreal offre \textbf{Set Timer by Event} (e \emph{by Function Name}):
si collega un evento personalizzato, si imposta un tempo (per esempio 60
secondi), e l'engine chiama quell'evento allo scadere, senza far nulla a ogni
frame. Con \emph{Get Timer Remaining Time by Handle} si legge quanto manca, per
mostrarlo a schermo. La somma dei Delta Seconds resta comunque il modo più
trasparente per capire che cosa sta succedendo.

% ---------------------------------------------------------------------
\section{Il trampolino: Jump e Launch Character}
\label{gd12:sec:trampolino}

Francesco vuole una \textbf{piattaforma su cui si atterra e si salta
automaticamente}. ``Chiaro, chiaro.''
\begin{enumerate}
  \item ``Fai la stessa cosa che hai fatto sulla lava'': servono le
        informazioni del giocatore, quindi overlap e \textbf{cast} al
        personaggio (``puoi anche fare Ctrl+C, Ctrl+V'').
  \item ``Devi collegare chiaramente anche il filo bianco, se no non va il
        codice.''
  \item Dal riferimento \emph{As ThirdPersonCharacter}, il nodo
        \textbf{Jump}: ``fai saltare il player; vedi, il target è il
        ThirdPersonCharacter, che è il player che hai appena toccato''.
\end{enumerate}

\paragraph{Quanto salta?} Con \textbf{Jump} il personaggio salta ``come se tu
stessi premendo salto'', con quella forza lì. Per un'altra forza c'è
\textbf{Launch Character}, sempre tirato fuori dal riferimento al personaggio:
gli si dà una \textbf{velocità di lancio}. ``Considera che è un launch character,
quindi può anche andare in orizzontale, non solo in verticale'': il vettore ha
$x$, $y$ e $z$.

\begin{figure}[htbp]
\centering
\begin{bpgraph}
  \bpnode[2]{bo}{0,0}{bpHEvent}{On Component Begin Overlap\\\bpsub{(Box)}}{4.0}{}{exec/, obj/Other Actor}
  \bpnode{ca}{4.9,0}{bpHCast}{Cast To ThirdPersonCharacter}{4.6}{exec/, obj/Object}%
    {exec/, exec/Cast Failed, obj/As BP Third Person Character}
  \bpnode[2]{lc}{10.6,0}{bpHFunc}{Launch Character\\\bpsub{Target is Character}}{4.4}%
    {exec/, obj/Target, vec/Launch Velocity \bpval{0 \ 0 \ 1500}, bool/XYOverride, bool/ZOverride}{exec/}
  \bpwire{exec}{bo-o1}{ca-i1}
  \bpwire{obj}{bo-o2}{ca-i2}
  \bpwire{exec}{ca-o1}{lc-i1}
  \bpwire{obj}{ca-o3}{lc-i2}
  \bpnote{4.9,-2.15}{con \emph{Jump} al posto di Launch Character\\il salto è quello normale del personaggio}
\end{bpgraph}
\caption{Il trampolino: quando il personaggio entra nella box viene lanciato
verso l'alto. Il valore va cercato provando: 20 o 200 sono troppo poco per
accorgersene.}
\label{gd12:fig:trampolino}
\end{figure}

\paragraph{Il debug passo per passo.} All'inizio non succede niente, ed è
l'occasione per un metodo.
\begin{itemize}
  \item \textbf{Il valore.} ``20 è veramente poco. 2 è tipo niente.'' Conviene
        lavorare sulla \textbf{$z$}, ``che è quella che ti lancia'' in alto;
        lo studente aveva modificato un asse orizzontale.
  \item \textbf{La box.} Nella viewport, l'area va allargata e alzata un
        po', così che il personaggio la tocchi davvero.
  \item \textbf{Test rapidi.} ``Quando fai test rapidi, non fare tutto quanto
        il livello e poi mettere quello lì in fondo, sennò diventa lungo tutte
        quante le volte.'' Si sposta il \textbf{Player Start} vicino a ciò che
        si sta testando (``mettilo dove ti viene più comodo'').
  \item \textbf{Il Print String.} Uno \textbf{Print String} proprio dopo il
        Begin Overlap dice se il codice \emph{arriva} lì: ``così capiamo dov'è
        il problema''. Il messaggio compare: il codice arriva.
\end{itemize}
Il problema, alla fine, era solo il numero: con 200 il lancio è impercettibile;
con ``un numero molto più alto'', ``direi che funziona''. ``Funzionava tutto,
semplicemente il valore era troppo basso e non te ne accorgevi.''

% ---------------------------------------------------------------------
\section{Il giro tra i banchi, seconda parte}
\label{gd12:sec:giro2}

\paragraph{Il mulino e il vento.} Antonino condivide di nuovo: nella sua scena
c'è una pala che gira (il movimento è del modello, il codice ce l'ha messo lui).
``Puoi anche fare così, va bene. Fai attenzione che c'è il problema del delta
time, che rimane sempre'': la rotazione sull'Event Tick va moltiplicata per i
Delta Seconds. C'è anche una zona di \textbf{vento}, fatta con una box,
``un pochettino strano, però te lo concedo''. ``Direi che sei sulla retta via:
continua tranquillamente a implementare feature così.''

\paragraph{Il personaggio scalato.} Uno studente ha ingrandito molto il
personaggio, che ora si comporta in modo strano. ``Perché questo personaggio non
è pensato per essere scalato'': le \textbf{animazioni} non sono pensate per
quella scala, ``e quindi il codice non funziona correttamente''. Bisognerebbe
modificare l'animazione.

\paragraph{Una piattaforma che ruota.} Un'ultima studentessa fatica con la
timeline, che ``va avanti e indietro'' in modo strano. Il docente torna al
principio: la timeline non fa nulla da sola, ``semplicemente ti disegni questo
grafico; questa curva la disegni, e poi il valore di questa curva ti esce, e devi
sfruttarlo per far andare avanti e indietro''. Se la curva resta a zero, non
cambia niente; se sale e scende, l'oggetto sale e scende. Rifanno insieme i
punti: tempo 0 valore 0, tempo 2,5 valore 200, tempo 5 valore 0, ``così fa su e
poi giù''. Due consigli:
\begin{itemize}
  \item per spostare un oggetto partendo da dove si trova, prendere la
        \textbf{posizione attuale} (o tenere la base a zero e muovere un
        componente figlio), altrimenti, con valori che vanno sotto lo zero,
        ``poi è un casino'';
  \item per farlo \textbf{ruotare} si usa lo stesso principio, con un Set
        Relative Rotation; e se servono più grandezze insieme si aggiungono
        \textbf{più tracce} nella stessa timeline, ``così hai diversi
        grafici''.
\end{itemize}

\paragraph{Mandate i progressi.} Il docente ricorda a tutti: ``mandatemi poi il
video sulla chat mentre lavorate'', così sa come procede il lavoro. E a chi chiede
se si possano cambiare i parametri del personaggio: ``tutto quanto ciò che non è
specificato nel task potete modificarlo''.

% ---------------------------------------------------------------------
\section{La chiusura}
\label{gd12:sec:chiusura}

``Effettivamente sono le 17. Quindi direi che possiamo andarcene tutti quanti a
casa.'' Ultime domande? Uno studente riassume il task, e il docente conferma:
\begin{itemize}
  \item \textbf{tre livelli};
  \item dei \textbf{punti che possono far ripartire il livello} (trappole,
        cadute, il timer);
  \item il resto, ``se volete, sì; se no, niente''.
\end{itemize}
``Chiaramente, più robe fate, meglio le fate, migliore sarà il voto. Questo mi
sembra ovvio.''

Ci si rivede \textbf{lunedì} (la lezione su Git). Chi lavora già nel weekend può
scrivere i propri dubbi e fare domande. ``Allora, grazie mille. Ci
rivediamo.''

% ---------------------------------------------------------------------
\section*{Punti chiave del capitolo}
\addcontentsline{toc}{section}{Punti chiave del capitolo}

\begin{itemize}
  \item \textbf{Task} (consegna il \textbf{25}, zip su Google Drive): un
        \textbf{platformer 3D} dal template in terza persona, con \textbf{tre
        livelli}, \textbf{ostacoli} che possono far ricominciare il livello,
        \textbf{aree che modificano il movimento} e \textbf{un livello a
        tempo} che riparte se il timer scade. Asset da Fab ammessi e da
        segnalare con un link; codice tutto proprio; aggiornamenti informali
        nella chat della classe, come in un team (non valutati). Dal task
        successivo, consegne su GitHub.
  \item Compilare Unreal \textbf{dal sorgente} serve per modificare l'engine,
        buildare i \textbf{server dedicati}, buildare per \textbf{console} e
        provare le versioni in anteprima (Unreal Engine 6); costa spazio e
        un giorno di build. \textbf{Python} si usa per tool dell'editor, non per
        il codice di gioco.
  \item Un \textbf{array} contiene più valori dello stesso tipo in una sola
        variabile; gli \textbf{indici partono da 0}. \textbf{Add} aggiunge in
        fondo, \textbf{Remove Index} toglie una posizione. Mai tante variabili
        per la stessa cosa in posizioni diverse.
  \item Il \textbf{For Each Loop} esegue il Loop Body per ogni elemento, con
        Array Element e Array Index, e Completed alla fine.
        \textbf{Get All Actors Of Class} restituisce un array di tutti gli
        oggetti di una classe: ``prendi tutti, controlla uno per uno
        (distanza, Branch), agisci'' è alla base di moltissimo codice.
  \item Si consegna con \textbf{File $\rightarrow$ Zip Project}, non
        zippando la cartella a mano: \texttt{Config} e \texttt{Content} sono
        ciò che conta, le cartelle di cache si rigenerano, \texttt{Saved}
        contiene log e salvataggi.
  \item In \textbf{Project Settings $\rightarrow$ Maps \& Modes}: \emph{Editor
        Startup Map} (per lo sviluppatore) e \emph{Game Default Map} (per il
        giocatore, di solito menu o splash screen).
  \item \textbf{Packaging} da \emph{Platforms}: \textbf{Shipping} per la
        versione pubblica, \textbf{DebugGame} per console, log e test. Lento la
        prima volta. Non richiesto per il task.
  \item Nuovi livelli con \textbf{File $\rightarrow$ New Level $\rightarrow$
        Basic}, poi duplicati con Ctrl+C, Ctrl+V. \textbf{Fab}: asset, plugin,
        progetti completi; filtrare per formato Unreal Engine e gratuito; tre
        pacchetti gratuiti a tempo ogni quindici giorni; gli asset vanno
        \textbf{solo avanti di versione, mai indietro}.
  \item La \textbf{timeline} fa variare valori nel tempo: \emph{Float Track},
        chiavi con \emph{Add Key} (meglio scrivere tempo e valore a mano),
        \emph{Length}, \emph{Loop}, interpolazione \emph{Auto}; uscite
        \textbf{Update} (ogni frame), \textbf{Finished}, \textbf{Direction} e
        un pin per traccia. \textbf{Play} riprende da dove si trova,
        \textbf{Play from Start} riparte, \textbf{Reverse} torna indietro: con
        Begin e End Overlap si fa una porta. Per un andirivieni in loop, la
        curva deve tornare al valore di partenza.
  \item \textbf{Relative} è rispetto all'Actor che contiene il componente,
        \textbf{World} rispetto al mondo. Una location è solo posizione, una
        transform include rotazione e scala.
  \item Muovere sull'\textbf{Event Tick}: posizione attuale + spostamento
        $\times$ \textbf{Delta Seconds}. È il modo base, ad esempio per i
        proiettili; la timeline ci si appoggia sopra.
  \item \textbf{Get Current Level Name} $\rightarrow$ \textbf{Open Level (by
        Name)} ricarica il livello corrente: è la base di trappole e timer.
        Sempre un \textbf{cast al personaggio} prima, o l'overlap scatterà per
        qualsiasi oggetto (anche all'infinito).
  \item La \textbf{doppia porta}: una timeline, un booleano che fa da lucchetto,
        \emph{IsInFront} per il verso, Direction per decidere se richiudere.
        Le \textbf{Scene Component} fanno da cerniere, come un pivot.
  \item Un \textbf{timer} si costruisce sommando i Delta Seconds (sommare 1
        conta i frame). \textbf{Jump} fa saltare come il tasto,
        \textbf{Launch Character} lancia con un vettore anche orizzontale; per
        il debug: Player Start vicino, Print String dopo l'evento, valori
        abbastanza grandi.
  \item \textbf{Engine Scalability $\rightarrow$ Low} per i PC che faticano;
        il personaggio del template non va scalato troppo, perché le
        animazioni non lo prevedono.
\end{itemize}
